\documentclass[11pt,a4paper,openany]{report}

% ------------------------------------------------------------------
% Idioma, fuentes y página
% ------------------------------------------------------------------
\usepackage[T1]{fontenc}
\usepackage[utf8]{inputenc}
\usepackage[english]{babel}
\babelprovide[import,main]{spanish}
\usepackage{lmodern}
\usepackage{microtype}
\usepackage[a4paper,margin=2.4cm,headheight=14pt]{geometry}
\usepackage{fancyhdr}
\usepackage{amsmath,amssymb,bm}
\usepackage{siunitx}
\sisetup{output-decimal-marker={,},group-separator={\,},per-mode=symbol}
\usepackage{booktabs,array,tabularx,multirow}
\usepackage{enumitem}
\usepackage{float}
\usepackage{caption,subcaption}
\captionsetup{font=small,labelfont=bf}
\usepackage{xcolor}

% ------------------------------------------------------------------
% Gráficos
% ------------------------------------------------------------------
\usepackage{tikz}
\usetikzlibrary{arrows.meta,positioning,calc,patterns,shapes.geometric,
  decorations.pathmorphing,decorations.markings,fit,backgrounds,matrix,shadows}
\usepackage{pgfplots}
\pgfplotsset{compat=1.18}
\usepgfplotslibrary{groupplots}
\pgfplotsset{
  manual/.style={width=\linewidth,height=4.6cm,grid=major,
    grid style={gray!25},tick label style={font=\footnotesize},
    label style={font=\small},legend style={font=\footnotesize,
    fill=white,fill opacity=0.85,text opacity=1,draw=gray!50},
    every axis plot/.append style={line width=0.9pt},
    /pgf/number format/use comma},
}
\definecolor{c1}{RGB}{0,114,189}
\definecolor{c2}{RGB}{217,83,25}
\definecolor{c3}{RGB}{237,177,32}
\definecolor{c4}{RGB}{126,47,142}
\definecolor{c5}{RGB}{119,172,48}
\definecolor{c6}{RGB}{77,190,238}

% ------------------------------------------------------------------
% Código
% ------------------------------------------------------------------
\usepackage{listings}
\lstdefinestyle{matlab}{
  language=Matlab, basicstyle=\ttfamily\footnotesize,
  keywordstyle=\color{blue!70!black}\bfseries,
  commentstyle=\color{green!45!black}\itshape,
  stringstyle=\color{violet},
  numbers=left, numberstyle=\tiny\color{gray}, numbersep=6pt,
  frame=single, rulecolor=\color{gray!40}, backgroundcolor=\color{gray!5},
  breaklines=true, columns=fullflexible, keepspaces=true,
  showstringspaces=false, xleftmargin=1.6em, framexleftmargin=1.4em,
  morekeywords={classdef,properties,methods,persistent,assert},
  literate={á}{{\'a}}1 {é}{{\'e}}1 {í}{{\'i}}1 {ó}{{\'o}}1 {ú}{{\'u}}1
           {ñ}{{\~n}}1 {Á}{{\'A}}1 {É}{{\'E}}1 {Ó}{{\'O}}1 {ü}{{\"u}}1
           {θ}{{$\theta$}}1 {°}{{$^\circ$}}1 {±}{{$\pm$}}1 {→}{{$\to$}}1
           {─}{{-}}1 {►}{{$\triangleright$}}1 {┐}{{+}}1 {┘}{{+}}1 {┼}{{+}}1
           {│}{{|}}1 {◄}{{$\triangleleft$}}1 {┴}{{+}}1
}
\lstset{style=matlab}

% ------------------------------------------------------------------
% Cajas
% ------------------------------------------------------------------
\usepackage[most]{tcolorbox}
\newtcolorbox{nota}[1][]{colback=c1!5,colframe=c1!70,fonttitle=\bfseries,
  title={Nota},breakable,#1}
\newtcolorbox{aviso}[1][]{colback=c2!6,colframe=c2!80,fonttitle=\bfseries,
  title={Atención},breakable,#1}
\newcounter{ejemplo}[chapter]
\renewcommand{\theejemplo}{\thechapter.\arabic{ejemplo}}
\newtcolorbox{ejemplo}[1]{colback=c5!6,colframe=c5!60!black,fonttitle=\bfseries,
  breakable,before upper={\refstepcounter{ejemplo}},
  title={Ejemplo \stepcounter{ejemplo}\theejemplo\addtocounter{ejemplo}{-1}: #1}}
\newcounter{practica}
\newtcolorbox{practica}[1]{enhanced,breakable,colback=white,colframe=c4!70!black,
  fonttitle=\bfseries\large,title={Práctica \thepractica: #1},
  before upper={},attach boxed title to top left={xshift=4mm,yshift=-2mm},
  boxed title style={colback=c4!70!black}}
\newcommand{\objetivos}{\par\smallskip\noindent\textbf{Objetivos}\par}
\newcommand{\procedimiento}{\par\smallskip\noindent\textbf{Procedimiento}\par}
\newcommand{\preguntas}{\par\smallskip\noindent\textbf{Cuestiones}\par}
\newcommand{\entregables}{\par\smallskip\noindent\textbf{Entregables}\par}

\usepackage[hidelinks]{hyperref}
\hypersetup{pdftitle={Manual del péndulo lineal (carro-péndulo)},
  pdfauthor={Serie de experimentos de sistemas de control}}

% ------------------------------------------------------------------
% Macros
% ------------------------------------------------------------------
\newcommand{\dd}{\mathrm{d}}
\newcommand{\vect}[1]{\bm{#1}}
\newcommand{\code}[1]{\texttt{#1}}
\newcommand{\xd}{\dot{x}}
\newcommand{\xdd}{\ddot{x}}
\newcommand{\thd}{\dot{\theta}}
\newcommand{\thdd}{\ddot{\theta}}
\newcommand{\Jp}{J_p}
\renewcommand{\figurename}{Figura}
\renewcommand{\tablename}{Tabla}
\renewcommand{\contentsname}{Índice}
\renewcommand{\listfigurename}{Índice de figuras}
\renewcommand{\chaptername}{Capítulo}
\renewcommand{\appendixname}{Apéndice}
\renewcommand{\bibname}{Bibliografía}
\renewcommand{\lstlistingname}{Listado}

\pagestyle{fancy}
\fancyhf{}
\fancyhead[L]{\small\nouppercase{\leftmark}}
\fancyhead[R]{\small Péndulo lineal}
\fancyfoot[C]{\thepage}
\renewcommand{\headrulewidth}{0.4pt}
\fancypagestyle{plain}{\fancyhf{}\fancyfoot[C]{\thepage}\renewcommand{\headrulewidth}{0pt}}

\setlength{\parskip}{0.35em}
\setlist{itemsep=0.15em,topsep=0.3em}

\begin{document}

% ------------------------------------------------------------------
% Portada
% ------------------------------------------------------------------
\begin{titlepage}
\centering
\vspace*{1.2cm}
{\large Serie de experimentos de sistemas de control\par}
\vspace{0.4cm}
{\color{c1}\rule{\linewidth}{1.2pt}\par}
\vspace{0.4cm}
{\Huge\bfseries Péndulo lineal\\[0.2em] (carro-péndulo)\par}
\vspace{0.4cm}
{\Large Manual teórico y práctico\par}
\vspace{0.2cm}
{\large Modelado, análisis, control y simulación en MATLAB/Simulink\par}
{\color{c1}\rule{\linewidth}{1.2pt}\par}
\vspace{1.2cm}
\begin{tikzpicture}[scale=1.25]
  \fill[pattern=north east lines,pattern color=gray!60] (-3.6,-0.55) rectangle (3.6,-0.4);
  \draw[very thick] (-3.6,-0.4) -- (3.6,-0.4);
  \draw[line width=3pt] (-3.4,-0.4) -- (-3.4,0.1);
  \draw[line width=3pt] (3.4,-0.4) -- (3.4,0.1);
  \draw[fill=c1!70,rounded corners=2pt] (-0.75,-0.35) rectangle (0.75,0.35);
  \fill (-0.45,-0.4) circle (0.07); \fill (0.45,-0.4) circle (0.07);
  \draw[line width=4pt,c2,line cap=round] (0,0) -- ++(70:2.6);
  \fill[c2] (0,0) ++(70:2.6) circle (0.12);
  \draw[fill=white] (0,0) circle (0.07);
  \draw[-{Latex[length=3mm]},line width=1.5pt,c4] (-1.9,0) -- (-0.8,0) node[midway,above]{$F$};
  \draw[dashed,gray] (0,0) -- (0,2.7);
  \draw[->] (0,1.1) arc (90:70:1.1) node[midway,above]{$\theta$};
  \draw[-{Latex},thick] (-3.0,-0.9) -- (-2.0,-0.9) node[right]{$x$};
\end{tikzpicture}
\vfill
{\large Versión 1.0 \quad·\quad Noviembre 2026\par}
\vspace{0.6cm}
{\small Acompaña al paquete de software \code{linear\_pendulum} (\code{cp\_build\_model}, \code{CartPole}, \code{cp\_gui}).\par}
\end{titlepage}

\tableofcontents
\listoffigures

\chapter{Introducción}
\label{cap:intro}

El \emph{péndulo invertido sobre carro} (carro-péndulo o péndulo lineal) es el banco de pruebas clásico de la teoría de control. Reúne en un sistema pequeño casi todas las dificultades que aparecen en la práctica: un equilibrio inestable, una no linealidad fuerte, un sistema subactuado (dos grados de libertad y una sola entrada), un cero de fase no mínima, saturación del actuador, límites físicos de recorrido y sensores que solo miden posiciones.

Este manual acompaña al paquete de simulación \code{linear\_pendulum} y está organizado en dos partes:

\begin{itemize}
  \item \textbf{Parte teórica} (capítulos~\ref{cap:modelo}--\ref{cap:control}). Se deriva el modelo, se linealiza, se analizan sus propiedades (estabilidad, controlabilidad, observabilidad, fase no mínima), se estudia la instrumentación y se desarrollan cuatro estrategias de control: LQR, PID+PD, \emph{swing-up} por energía y grúa con anti-balanceo.
  \item \textbf{Parte práctica} (capítulos~\ref{cap:simulink}--\ref{cap:practicas}). Se describe el modelo de Simulink, la interfaz de código \code{CartPole} y la interfaz gráfica \code{cp\_gui}. El manual termina con diez prácticas guiadas.
\end{itemize}

Todas las figuras de resultados de este manual se han obtenido simulando el modelo no lineal completo, con actuador, cuantización de los encoders y controlador discreto, y usando los parámetros por defecto de \code{cp\_params.m}. Por tanto, deberían reproducirse con el software tal y como se entrega.

\section{El sistema físico}

La figura~\ref{fig:sistema} muestra el montaje. Un carro de masa $M$ se desplaza sobre un raíl horizontal de longitud útil $2x_{\lim}=\SI{1.6}{m}$. Un actuador lineal (por ejemplo, un motor con correa o un motor lineal) ejerce sobre él una fuerza $F$. En el carro va articulada una varilla uniforme de masa $m$ y longitud $L_p$, que gira libremente en el plano vertical. Dos encoders miden la posición del carro $x$ y el ángulo de la varilla $\theta$.

\begin{figure}[H]
\centering
\begin{tikzpicture}[scale=1.5,>=Latex]
  % raíl
  \fill[pattern=north east lines,pattern color=gray!60] (-4,-0.62) rectangle (4,-0.5);
  \draw[thick] (-4,-0.5) -- (4,-0.5);
  \draw[line width=3pt] (-3.7,-0.5) -- (-3.7,0.1) node[above,font=\small]{tope};
  \draw[line width=3pt] (3.7,-0.5) -- (3.7,0.1) node[above,font=\small]{tope};
  % carro
  \draw[fill=c1!60,rounded corners=2pt] (-0.7,-0.4) rectangle (0.7,0.3);
  \node[font=\small] at (0,-0.1) {$M$};
  \draw[fill=gray] (-0.42,-0.45) circle (0.06); \draw[fill=gray] (0.42,-0.45) circle (0.06);
  % varilla
  \def\ang{25}
  \draw[line width=5pt,c2!80,line cap=round] (0,0.15) -- ++(90-\ang:3.0) coordinate (tip);
  \coordinate (piv) at (0,0.15);
  \path (piv) -- (tip) coordinate[pos=0.5] (cm);
  \fill[black] (cm) circle (0.07);
  \node[right=3pt,font=\small] at (cm) {CdM, $m$, $J$};
  \draw[fill=white] (piv) circle (0.06);
  % ejes y ángulo
  \draw[dashed,gray] (piv) -- ++(0,3.1) node[above,black,font=\small]{vertical};
  \draw[->,thick] (piv) ++(0,1.3) arc (90:90-\ang:1.3);
  \node[font=\small] at ($(piv)+(0.28,1.45)$) {$\theta$};
  % cotas l y Lp
  \draw[<->,gray!80!black] ($(piv)+(-0.25,0)+(180-\ang:0)$) -- ++(90-\ang:1.5) node[midway,left,font=\small]{$l$};
  \draw[<->] ($(piv)+(0.35,-0.15)$) ++(90-\ang:0) -- ++(90-\ang:3.0);
  \node[font=\small,right] at ($(piv)!0.75!(tip)+(0.35,-0.1)$) {$L_p=2l$};
  % fuerzas
  \draw[->,line width=1.6pt,c4] (-2.0,-0.05) -- (-0.75,-0.05) node[midway,above]{$F$};
  \draw[->,line width=1.4pt,c5!70!black] ($(tip)+(-1.1,0)$) -- ($(tip)+(-0.12,0)$) node[pos=0.3,above]{$F_d$};
  \draw[->,thick,gray!70!black] ($(cm)+(0.9,0)$) -- ++(0,-0.7) node[right,font=\small]{$m g$};
  % coordenada x
  \draw[->,thick] (-3.6,-1.0) -- (3.6,-1.0) node[right]{$x$};
  \draw (0,-0.95) -- (0,-1.05) node[below,font=\small]{$x$};
  \draw (-3.7,-0.95) -- (-3.7,-1.05) node[below,font=\small]{$-x_{\lim}$};
  \draw (3.7,-0.95) -- (3.7,-1.05) node[below,font=\small]{$x_{\lim}$};
  % encoders
  \node[draw,circle,fill=yellow!40,inner sep=1pt,font=\scriptsize] at (0.45,0.05) {E$_\theta$};
  \node[draw,rectangle,fill=yellow!40,inner sep=2pt,font=\scriptsize,anchor=north] at (-2.6,-0.62) {Encoder lineal E$_x$};
\end{tikzpicture}
\caption{Esquema del péndulo lineal y convención de signos. $\theta=0$ corresponde a la varilla vertical hacia arriba; $\theta>0$ indica inclinación hacia $+x$.}
\label{fig:sistema}
\end{figure}

\section{Convenciones}

\begin{table}[H]
\centering\small
\begin{tabular}{@{}lll@{}}
\toprule
Símbolo & Significado & Unidad\\
\midrule
$x$ & posición del carro (positiva a la derecha) & m\\
$\theta$ & ángulo de la varilla respecto a la vertical superior; $\theta=\pi$ colgando & rad\\
$F$ & fuerza aplicada por el actuador sobre el carro & N\\
$F_d$ & fuerza de perturbación horizontal en la punta de la varilla & N\\
$\vect{s}$ & estado $[x,\ \xd,\ \theta,\ \thd]^\top$ & --\\
$\vect{y}$ & medida $[x_m,\ \theta_m]^\top$ (solo posiciones) & m, rad\\
$r$ & referencia de posición del carro & m\\
\bottomrule
\end{tabular}
\caption{Convenciones de signos y variables.}
\end{table}

\section{Contenido del paquete de software}

\begin{table}[H]
\centering\small
\begin{tabularx}{\linewidth}{@{}lX@{}}
\toprule
Archivo & Función\\
\midrule
\code{cp\_params.m} & Todos los parámetros físicos, de sensores, temporales y de los controladores.\\
\code{cp\_linearize.m} & Matrices $(A,B,C,D)$ en torno a $\theta_0=0$ o $\theta_0=\pi$.\\
\code{cp\_lqr.m}, \code{cp\_design.m} & Cálculo de las ganancias LQR (con o sin Control System Toolbox).\\
\code{cp\_plant\_block.m} & Dinámica no lineal (código del bloque \emph{MATLAB Function} de la planta).\\
\code{cp\_controller\_block.m} & Controlador discreto con 5 modos (código del bloque del controlador).\\
\code{cp\_ref\_block.m} & Generador de referencias.\\
\code{cp\_build\_model.m} & Genera el modelo \code{cartpole\_sim.slx} mediante un script.\\
\code{CartPole.m} & Clase con la interfaz de código: ensayos por lotes y control en vivo.\\
\code{cp\_gui.m} & Interfaz gráfica en vivo.\\
\code{cp\_draw.m}, \code{cp\_animate.m}, \code{cp\_plot\_results.m} & Animación y gráficas.\\
\code{examples.m} & Recorrido guiado por secciones.\\
\code{python\_reference\_sim.py} & Réplica independiente en Python usada para validar el modelo.\\
\code{validation/} & Scripts de validación cruzada Simulink--Python (apéndice~\ref{cap:validacion}).\\
\bottomrule
\end{tabularx}
\caption{Archivos del paquete \code{linear\_pendulum}.}
\label{tab:archivos}
\end{table}

\section{Requisitos y puesta en marcha}

Se necesita MATLAB y Simulink R2020b o posterior. La Control System Toolbox es opcional: si no está instalada, \code{cp\_lqr} resuelve la ecuación de Riccati por el método del Hamiltoniano (sección~\ref{sec:hamiltoniano}).

\begin{lstlisting}[caption={Primeros pasos.}]
cd src/linear_pendulum
cp_build_model                 % crea cartpole_sim.slx (solo la primera vez)
cp = CartPole();               % interfaz de código
res = cp.run(6);               % ensayo LQR por defecto (theta0 = 0.1 rad)
cp_plot_results(res); cp.animate(res);
cp_gui(cp);                    % interfaz gráfica en vivo
\end{lstlisting}

\chapter{Modelado dinámico}
\label{cap:modelo}

\section{Coordenadas generalizadas}

El sistema tiene dos grados de libertad, que tomamos como coordenadas generalizadas: $\vect{q}=[x,\ \theta]^\top$. La varilla se modela como un sólido rígido. Su centro de masas (CdM) está a una distancia $l=L_p/2$ del pivote, y su momento de inercia respecto al CdM es $J=mL_p^2/12$. La posición del CdM es
\begin{equation}
  x_G = x + l\sin\theta, \qquad y_G = l\cos\theta,
\end{equation}
y su velocidad
\begin{equation}
  \dot{x}_G = \xd + l\cos\theta\,\thd, \qquad \dot{y}_G = -l\sin\theta\,\thd .
\end{equation}

\section{Energías y lagrangiano}

La energía cinética es la suma de la del carro, la de traslación del CdM de la varilla y la de rotación de la varilla:
\begin{align}
  T &= \tfrac12 M\xd^2 + \tfrac12 m\left(\dot{x}_G^2+\dot{y}_G^2\right) + \tfrac12 J\thd^2 \nonumber\\
    &= \tfrac12 (M+m)\xd^2 + m l\cos\theta\,\xd\thd + \tfrac12\left(J+ml^2\right)\thd^2 .
  \label{eq:T}
\end{align}
Llamaremos $\Jp = J+ml^2$ al momento de inercia de la varilla respecto al pivote (para la varilla uniforme, $\Jp=mL_p^2/3$). Tomando la altura del pivote como referencia, la energía potencial es
\begin{equation}
  V = m g l\cos\theta ,
\end{equation}
y el lagrangiano $\mathcal{L}=T-V$.

\section{Ecuaciones de Euler--Lagrange}

Las fuerzas generalizadas no conservativas son la fuerza del actuador, las fricciones, los topes del raíl y la perturbación $F_d$, que se aplica horizontalmente en la punta de la varilla, en $(x+L_p\sin\theta,\ L_p\cos\theta)$. El trabajo virtual de $F_d$ vale $F_d\,\delta(x+L_p\sin\theta)=F_d\,\delta x + F_dL_p\cos\theta\,\delta\theta$, de modo que
\begin{equation}
  Q_x = F + F_d - F_f(\xd) + F_w(x,\xd), \qquad Q_\theta = -b_p\thd + F_dL_p\cos\theta .
\end{equation}
Aplicando $\frac{\dd}{\dd t}\frac{\partial\mathcal L}{\partial\dot q_i}-\frac{\partial\mathcal L}{\partial q_i}=Q_i$:

\paragraph{Coordenada $x$.} $\frac{\partial\mathcal L}{\partial\xd}=(M+m)\xd+ml\cos\theta\,\thd$, y $\frac{\partial\mathcal L}{\partial x}=0$, de donde
\begin{equation}
  (M+m)\xdd + ml\cos\theta\,\thdd - ml\sin\theta\,\thd^2 = F + F_d - F_f + F_w .
  \label{eq:el1}
\end{equation}

\paragraph{Coordenada $\theta$.} $\frac{\partial\mathcal L}{\partial\thd}=ml\cos\theta\,\xd+\Jp\thd$, y $\frac{\partial\mathcal L}{\partial\theta}=-ml\sin\theta\,\xd\thd+mgl\sin\theta$. Por tanto,
\begin{equation}
  ml\cos\theta\,\xdd + \Jp\thdd - mgl\sin\theta = -b_p\thd + F_dL_p\cos\theta .
  \label{eq:el2}
\end{equation}

\section{Forma matricial}

Las ecuaciones~\eqref{eq:el1}--\eqref{eq:el2} se escriben como
\begin{equation}
\underbrace{\begin{bmatrix} M+m & ml\cos\theta\\ ml\cos\theta & \Jp\end{bmatrix}}_{\mathbf{M}(\theta)}
\begin{bmatrix}\xdd\\ \thdd\end{bmatrix}
=
\underbrace{\begin{bmatrix} F + F_d - F_f + F_w + ml\sin\theta\,\thd^2\\ mgl\sin\theta - b_p\thd + F_dL_p\cos\theta\end{bmatrix}}_{\vect{h}(\vect{q},\dot{\vect{q}},F,F_d)} .
\label{eq:matricial}
\end{equation}
La matriz de masas $\mathbf{M}(\theta)$ es simétrica y definida positiva para todo $\theta$, porque
\begin{equation}
  \det\mathbf{M} = (M+m)\Jp - m^2l^2\cos^2\theta \ \ge\ (M+m)\Jp-m^2l^2 = MJ_p + mJ > 0 .
\end{equation}
Así que siempre puede invertirse, y el modelo de estado no lineal queda
\begin{equation}
  \dot{\vect{s}} = \vect{f}(\vect{s},F,F_d) =
  \begin{bmatrix}\xd\\ [\mathbf{M}^{-1}\vect{h}]_1\\ \thd\\ [\mathbf{M}^{-1}\vect{h}]_2\end{bmatrix}.
  \label{eq:nl}
\end{equation}
Esto es exactamente lo que implementa \code{cp\_plant\_block.m}, que resuelve el sistema $2\times2$ con \code{Mt\textbackslash rhs} en cada evaluación.

\begin{ejemplo}{Matriz de masas con los parámetros por defecto}
Con $M=\SI{0.5}{kg}$, $m=\SI{0.2}{kg}$ y $L_p=\SI{0.6}{m}$ se obtiene $l=\SI{0.3}{m}$, $J=\SI{0.006}{kg.m^2}$ y $\Jp=\SI{0.024}{kg.m^2}$. En $\theta=0$:
\[
\mathbf{M}(0)=\begin{bmatrix}0{,}70 & 0{,}06\\ 0{,}06 & 0{,}024\end{bmatrix},\qquad
\det\mathbf{M}(0)=0{,}0168-0{,}0036=\SI{0.0132}{}.
\]
El término de acoplamiento $ml=0{,}06$ es grande comparado con $\Jp$. Por eso una aceleración del carro produce una aceleración angular importante, y eso es lo que permite controlar la varilla actuando solo sobre el carro.
\end{ejemplo}

\section{Fricción, topes y actuador}

\paragraph{Fricción del carro.} Se combinan una componente viscosa y una de Coulomb. La de Coulomb se suaviza con una tangente hiperbólica para evitar la discontinuidad en $\xd=0$:
\begin{equation}
  F_f(\xd) = b_c\,\xd + F_c\tanh\!\left(\xd/v_\varepsilon\right),
\end{equation}
con $b_c=\SI{0.5}{N.s/m}$, $F_c=\SI{0.2}{N}$ y $v_\varepsilon=\SI{0.01}{m/s}$. En el pivote solo hay fricción viscosa, $b_p=\SI{0.002}{N.m.s/rad}$.

\paragraph{Topes del raíl.} Si $|x|>x_{\lim}$, aparece una fuerza de tipo muelle--amortiguador:
\begin{equation}
  F_w = -k_w\big(x\mp x_{\lim}\big) - b_w\xd,\qquad k_w=\SI{1e4}{N/m},\ b_w=\SI{50}{N.s/m}.
\end{equation}

\paragraph{Actuador.} El actuador real (amplificador más motor) se modela como un sistema de primer orden seguido de saturación:
\begin{equation}
  \frac{F_{\text{act}}(s)}{F_{\text{cmd}}(s)} = \frac{1}{\tau_a s+1},\qquad \tau_a=\SI{10}{ms},\qquad |F_{\text{act}}|\le F_{\max}=\SI{15}{N}.
\end{equation}

\begin{nota}
Si se quiere modelar un motor de corriente continua con correa dentada, basta con sustituir este bloque por $F=\frac{k_t}{R_a r_p}\left(V-\frac{k_e}{r_p}\xd\right)$, donde $r_p$ es el radio de la polea. El término $-\frac{k_tk_e}{R_ar_p^2}\xd$ actúa como una fricción viscosa adicional (fuerza contraelectromotriz), así que el resto del manual sigue siendo válido cambiando $b_c$.
\end{nota}

\section{Puntos de equilibrio}

Con $F=F_d=0$, los equilibrios cumplen $\xd=\thd=0$ y $\sin\theta=0$:
\begin{itemize}
  \item $\theta=0$ (varilla arriba): equilibrio \textbf{inestable}. Es el objetivo de los modos LQR, PID y \emph{swing-up}.
  \item $\theta=\pi$ (varilla colgando): equilibrio \textbf{estable}, aunque poco amortiguado. Es el punto de trabajo del modo grúa.
\end{itemize}
En ambos casos, la posición $x$ es arbitraria: el sistema es invariante a traslaciones horizontales. Esto tiene consecuencias en el análisis lineal (un polo en el origen) y en el seguimiento de referencias (sección~\ref{sec:seguimiento}).

\section{Energía del péndulo}

La energía mecánica de la varilla \emph{respecto a un sistema de referencia que se mueve con el carro} se usará en el \emph{swing-up}. La definimos de modo que valga cero con la varilla en reposo arriba:
\begin{equation}
  E = \tfrac12\Jp\thd^2 + mgl(\cos\theta-1).
  \label{eq:energia}
\end{equation}
Colgando en reposo vale $E=-2mgl=\SI{-1.18}{J}$. Derivando y sustituyendo~\eqref{eq:el2} con $F_d=0$, queda
\begin{equation}
  \dot E = \Jp\thd\thdd - mgl\sin\theta\,\thd = -ml\cos\theta\,\thd\,\xdd - b_p\thd^2 .
  \label{eq:Edot}
\end{equation}
Es decir, \textbf{la aceleración del carro es la entrada que cambia la energía de la varilla}. Este resultado es la base de la ley de \emph{swing-up} de la sección~\ref{sec:swingup}.

\chapter{Linealización y análisis}
\label{cap:analisis}

\section{Linealización en torno a un equilibrio}

Sea $\theta_0\in\{0,\pi\}$ y $c_0=\cos\theta_0=\pm1$. Para desviaciones pequeñas $\varphi=\theta-\theta_0$ se cumple que $\sin\theta\approx c_0\varphi$, que $\cos\theta\approx c_0$ y que los términos $\thd^2$ son de segundo orden. Además se desprecia la fricción de Coulomb, que no es diferenciable en $\xd=0$. Con estas aproximaciones, el sistema~\eqref{eq:matricial} queda
\begin{equation}
\begin{bmatrix} M+m & mlc_0\\ mlc_0 & \Jp\end{bmatrix}
\begin{bmatrix}\xdd\\ \ddot\varphi\end{bmatrix}
=
\begin{bmatrix} F - b_c\xd\\ mglc_0\,\varphi - b_p\dot\varphi\end{bmatrix} .
\end{equation}
Con $D=(M+m)\Jp-m^2l^2$ (que no depende de $c_0$, porque $c_0^2=1$), al invertir se obtiene
\begin{align}
\xdd &= \frac{1}{D}\Big[\Jp(F-b_c\xd) - mlc_0\,(mglc_0\varphi - b_p\dot\varphi)\Big],\\
\ddot\varphi &= \frac{1}{D}\Big[-mlc_0(F-b_c\xd) + (M+m)(mglc_0\varphi - b_p\dot\varphi)\Big].
\end{align}
En forma de estado, con $\vect{s}=[x,\ \xd,\ \varphi,\ \dot\varphi]^\top$:
\begin{equation}
A=\frac1D\begin{bmatrix}
0 & D & 0 & 0\\
0 & -\Jp b_c & -m^2gl^2 & mlc_0b_p\\
0 & 0 & 0 & D\\
0 & mlc_0b_c & (M+m)mglc_0 & -(M+m)b_p
\end{bmatrix},\qquad
B=\frac1D\begin{bmatrix}0\\ \Jp\\ 0\\ -mlc_0\end{bmatrix},
\label{eq:AB}
\end{equation}
y como solo se miden posiciones,
\begin{equation}
C=\begin{bmatrix}1&0&0&0\\0&0&1&0\end{bmatrix},\qquad D_{yu}=\vect{0}.
\end{equation}

\begin{ejemplo}{Matrices numéricas (función \code{cp\_linearize})}
Con los parámetros por defecto, $D=\num{0.0132}$. \textbf{Arriba} ($\theta_0=0$):
\[
A_\uparrow=\begin{bmatrix}
0&1&0&0\\ 0&-0{,}909&-2{,}676&0{,}009\\ 0&0&0&1\\ 0&2{,}273&31{,}21&-0{,}106\end{bmatrix},\qquad
B_\uparrow=\begin{bmatrix}0\\1{,}818\\0\\-4{,}545\end{bmatrix}.
\]
\textbf{Abajo} ($\theta_0=\pi$) cambian de signo los términos con $c_0$:
\[
A_\downarrow=\begin{bmatrix}
0&1&0&0\\ 0&-0{,}909&-2{,}676&-0{,}009\\ 0&0&0&1\\ 0&-2{,}273&-31{,}21&-0{,}106\end{bmatrix},\qquad
B_\downarrow=\begin{bmatrix}0\\1{,}818\\0\\4{,}545\end{bmatrix}.
\]
Comprobación en MATLAB: \code{[A,B,C,D] = cp\_linearize(cp\_params(), 0)}.
\end{ejemplo}

\section{Polos en lazo abierto}

\begin{table}[H]
\centering\small
\begin{tabular}{@{}lll@{}}
\toprule
Equilibrio & Polos (rad/s) & Interpretación\\
\midrule
Arriba & $0,\ -0{,}711,\ -5{,}754,\ \mathbf{+5{,}450}$ & polo inestable: la caída tiene $\tau=\SI{0.18}{s}$\\
Abajo & $0,\ -0{,}717,\ -0{,}149\pm j\,5{,}573$ & oscilación a \SI{0.89}{Hz} con $\zeta=0{,}027$\\
\bottomrule
\end{tabular}
\caption{Polos del modelo lineal.}
\label{tab:polos}
\end{table}

\begin{itemize}
  \item El \textbf{polo en el origen} refleja la invariancia a traslaciones: si se desplaza el carro y se deja quieto, no vuelve. El polo en $-0{,}71$ procede de la fricción viscosa del carro ($\approx b_c/(M+m)$).
  \item Sin fricción, los polos del péndulo arriba serían $\pm\sqrt{(M+m)mgl/D}=\pm\SI{5.59}{rad/s}$. La constante de tiempo de la caída, $1/5{,}45\approx\SI{0.18}{s}$, fija el ancho de banda mínimo del lazo de control: el controlador tiene que ser bastante más rápido que eso.
  \item Abajo, la frecuencia natural ($\SI{5.57}{rad/s}$) es \emph{mayor} que la de la varilla con el pivote fijo, $\omega_0=\sqrt{mgl/\Jp}=\SI{4.95}{rad/s}$. El motivo es que el carro libre retrocede al oscilar la varilla, lo que acorta la longitud equivalente del péndulo.
\end{itemize}

\section{Funciones de transferencia y fase no mínima}
\label{sec:fnm}

Con el equilibrio de arriba y $F$ como entrada:
\begin{align}
G_x(s)=\frac{X(s)}{F(s)} &= \frac{1{,}818\,s^2+0{,}152\,s-44{,}59}{s\,(s^3+1{,}015\,s^2-31{,}14\,s-22{,}30)},\\[2pt]
G_\theta(s)=\frac{\Theta(s)}{F(s)} &= \frac{-4{,}545\,s}{s^3+1{,}015\,s^2-31{,}14\,s-22{,}30}.
\end{align}

\begin{figure}[H]
\centering
\begin{tikzpicture}
\begin{axis}[manual,width=0.75\linewidth,height=5.2cm,xlabel={$\mathrm{Re}(s)$},ylabel={$\mathrm{Im}(s)$},
  xmin=-7,xmax=7,ymin=-2,ymax=2,legend pos=north west,axis lines=middle,
  xlabel style={at={(ticklabel* cs:1)},anchor=north east},ylabel style={at={(ticklabel* cs:1)},anchor=south east}]
\addplot[only marks,mark=x,mark size=5pt,very thick,c2] coordinates {(0,0) (-0.711,0) (-5.754,0) (5.450,0)};
\addplot[only marks,mark=o,mark size=4pt,very thick,c1] coordinates {(-4.994,0) (4.911,0)};
\legend{polos,ceros de $G_x$}
\node[c2,above,font=\footnotesize] at (axis cs:5.45,0.15) {$+5{,}45$};
\node[c1,below,font=\footnotesize] at (axis cs:4.91,-0.2) {$+4{,}91$};
\end{axis}
\end{tikzpicture}
\caption{Mapa de polos y ceros del péndulo arriba (entrada $F$, salida $x$). $G_x$ tiene un cero en el semiplano derecho, muy cerca del polo inestable.}
\label{fig:pz}
\end{figure}

$G_x$ tiene un cero real positivo en $z=+4{,}91$. Sin fricción, ese cero estaría exactamente en $+\sqrt{mgl/\Jp}=\omega_0$. Esto tiene dos consecuencias:

\begin{enumerate}
  \item \textbf{Respuesta inversa.} Para llevar el carro a la derecha con la varilla en equilibrio, primero hay que inclinar la varilla hacia la derecha. Y para eso el carro tiene que moverse antes hacia la \emph{izquierda}. Cualquier controlador estabilizante lo hará, como se aprecia en la figura~\ref{fig:lqrstep}.
  \item \textbf{Límite de prestaciones.} Un cero en el semiplano derecho a frecuencia $z$ limita el ancho de banda alcanzable en $x$ a algo del orden de $z/2$. Además, el polo inestable en $p=5{,}45$ exige un ancho de banda mayor que $2p$ en el lazo del ángulo. Como $z$ y $p$ están tan cerca, el problema es intrínsecamente delicado; por eso el carro sigue con retraso las referencias rápidas.
\end{enumerate}

En $G_\theta$, el cero en el origen se cancela con el polo de traslación. Por eso, \textbf{midiendo solo $\theta$ no se puede saber dónde está el carro}, como confirma el análisis de observabilidad.

\section{Controlabilidad y observabilidad}

La matriz de controlabilidad $\mathcal{C}=[B\ AB\ A^2B\ A^3B]$ tiene rango~4, con número de condición $\kappa(\mathcal C)\approx 270$. Por tanto, el sistema es completamente controlable con una sola entrada $F$, aunque esté subactuado.

\begin{table}[H]
\centering\small
\begin{tabular}{@{}lcl@{}}
\toprule
Salidas medidas & $\operatorname{rango}\mathcal{O}$ & Conclusión\\
\midrule
$x$ y $\theta$ & 4 & observable (configuración del banco)\\
solo $x$ & 4 & observable: $\theta$ se deduce de cómo reacciona el carro\\
solo $\theta$ & 3 & \textbf{no observable}: la posición del carro queda indeterminada\\
\bottomrule
\end{tabular}
\caption{Observabilidad según los sensores disponibles.}
\end{table}

\begin{aviso}
Aunque bastaría con medir $x$, en la práctica se miden $x$ y $\theta$: estimar el ángulo a partir de $x$ es muy sensible a errores de modelo. Las velocidades $\xd$ y $\thd$ no se miden y hay que estimarlas, ya sea con derivadas filtradas (capítulo~\ref{cap:sensores}) o con un observador (práctica~10).
\end{aviso}

\section{Respuesta en frecuencia (equilibrio abajo)}

La figura~\ref{fig:bode} muestra el diagrama de Bode del modelo colgante, que es estable y por tanto identificable en lazo abierto (práctica~2). Destacan la resonancia a \SI{0.89}{Hz}, muy aguda (factor de calidad $1/2\zeta\approx19$, es decir, unos \SI{25}{dB}), y la antirresonancia de $G_x$ a la frecuencia $\omega_0/2\pi=\SI{0.79}{Hz}$. A esa frecuencia, al empujar el carro la varilla oscila de forma que «absorbe» el movimiento.

\begin{figure}[H]
\centering
\begin{tikzpicture}
\begin{groupplot}[group style={group size=1 by 2,vertical sep=0.5cm,xticklabels at=edge bottom},
  manual,height=4.2cm,xmode=log,xmin=0.05,xmax=10]
\nextgroupplot[ylabel={Magnitud [dB]},legend pos=south west]
\addplot[c1] table[x=f,y=mag]{data/bode_dn_x.dat};
\addplot[c2] table[x=f,y=mag]{data/bode_dn_th.dat};
\legend{$X/F$ [m/N],$\Theta/F$ [rad/N]}
\draw[dashed,gray] (axis cs:0.887,-120) -- (axis cs:0.887,40);
\draw[dashed,gray] (axis cs:0.788,-120) -- (axis cs:0.788,40);
\nextgroupplot[ylabel={Fase [$^\circ$]},xlabel={Frecuencia [Hz]}]
\addplot[c1] table[x=f,y=ph]{data/bode_dn_x.dat};
\addplot[c2] table[x=f,y=ph]{data/bode_dn_th.dat};
\end{groupplot}
\end{tikzpicture}
\caption{Bode del carro-péndulo colgante. Líneas discontinuas: antirresonancia de $X/F$ (\SI{0.79}{Hz}) y resonancia (\SI{0.89}{Hz}).}
\label{fig:bode}
\end{figure}

\chapter{Instrumentación y estimación de velocidades}
\label{cap:sensores}

\section{Cadena de medida}

En el modelo, cada medida pasa por tres etapas, que reproducen las de un banco real:
\begin{equation}
y_k = \mathcal{Q}_q\big(p(kT_s) + \sigma\,n_k\big),\qquad
\mathcal{Q}_q(u)=q\,\operatorname{round}(u/q),
\end{equation}
donde $p$ es la posición real ($x$ o $\theta$), $n_k\sim\mathcal N(0,1)$ es el ruido (desactivado por defecto) y $q$ es la resolución del encoder.

\begin{figure}[H]
\centering
\begin{tikzpicture}[node distance=1.0cm and 1.1cm,>=Latex,
  bloque/.style={draw,thick,rounded corners=2pt,minimum height=0.9cm,minimum width=1.6cm,align=center,font=\small,fill=c1!8},
  suma/.style={draw,circle,inner sep=1.5pt,thick}]
\node (p) {$p(t)$};
\node[suma,right=of p] (s) {$+$};
\node[bloque,right=of s] (q) {Cuantizador\\$q$};
\node[bloque,right=of q] (z) {ZOH\\$T_s$};
\node[right=of z] (y) {$y_k$};
\node[bloque,above=0.6cm of s,fill=gray!10] (n) {$\sigma\, n(t)$};
\draw[->] (p) -- (s); \draw[->] (n) -- (s); \draw[->] (s) -- (q); \draw[->] (q) -- (z); \draw[->] (z) -- (y);
\end{tikzpicture}
\caption{Modelo de cada canal de medida (subsistema \code{Sensores}).}
\end{figure}

\begin{table}[H]
\centering\small
\begin{tabular}{@{}lccc@{}}
\toprule
Canal & Resolución $q$ & Error de cuantización ($\sigma_q=q/\sqrt{12}$) & Ruido opcional $\sigma$\\
\midrule
$x$ & \SI{0.1}{mm} & \SI{0.029}{mm} & \SI{0.1}{mm}\\
$\theta$ & $2\pi/4096=\SI{1.53}{mrad}$ ($0{,}088^\circ$) & \SI{0.44}{mrad} & \SI{1}{mrad}\\
\bottomrule
\end{tabular}
\caption{Características de los encoders simulados.}
\end{table}

\section{Elección del periodo de muestreo}

El controlador trabaja con $T_s=\SI{2}{ms}$ (\SI{500}{Hz}). El polo más rápido del lazo cerrado LQR está en $|s|\approx\SI{14.6}{rad/s}$, es decir, unos \SI{2.3}{Hz}. La relación entre la frecuencia de muestreo y el ancho de banda resulta de unos 200, muy por encima del mínimo habitual (entre 20 y 30). Así, el controlador discreto se comporta prácticamente como el continuo, con el que se ha diseñado. El solver de la planta usa un paso fijo $h=\SI{1}{ms}=T_s/2$.

\section{Estimación de velocidades}

Derivar la señal cuantizada directamente amplifica mucho el ruido. Si $\theta$ cambia una sola cuenta entre dos muestras, la derivada salta en
\begin{equation}
\Delta\thd = q/T_s = \SI{1.53e-3}{}/\SI{0.002}{} \approx \SI{0.77}{rad/s} .
\end{equation}
Por eso se usa una \textbf{derivada filtrada}, equivalente en continuo a $\frac{s}{\tau_v s+1}$ con $\tau_v=\SI{10}{ms}$:
\begin{equation}
\hat v_k = a\,\hat v_{k-1} + (1-a)\,\frac{y_k-y_{k-1}}{T_s},\qquad a=e^{-T_s/\tau_v}=0{,}819.
\end{equation}
Su función de transferencia discreta es
\begin{equation}
\frac{\hat V(z)}{Y(z)} = \frac{1-a}{T_s}\,\frac{1-z^{-1}}{1-a z^{-1}} .
\end{equation}
Este filtro deriva correctamente hasta unos $1/\tau_v=\SI{100}{rad/s}$ (\SI{16}{Hz}) y atenúa por encima. El retraso que introduce a la frecuencia de cruce (unos \SI{15}{rad/s}) es de $\arctan(15\tau_v)\approx 8,5^\circ$, que el diseño tolera sin problemas.

\begin{figure}[H]
\centering
\begin{tikzpicture}
\begin{axis}[manual,height=5cm,xlabel={$t$ [s]},ylabel={$\thd$ [rad/s]},xmin=0,xmax=0.5,legend pos=south west,legend columns=3]
\addplot[gray!60,line width=0.5pt,const plot] table[x=t,y=raw]{data/velest.dat};
\addplot[c2,line width=1.2pt] table[x=t,y=filt]{data/velest.dat};
\addplot[black,dashed] table[x=t,y=true]{data/velest.dat};
\legend{diferencia directa,derivada filtrada,real}
\end{axis}
\end{tikzpicture}
\caption{Estimación de $\thd$ para $\theta(t)=0{,}2\sin(2\pi t)$, cuantizada a 4096 cuentas por vuelta y muestreada a \SI{500}{Hz}.}
\label{fig:velest}
\end{figure}

\begin{nota}
La práctica~9 propone variar $\tau_v$, $T_s$ y la resolución de los encoders, y observar cuándo el lazo empieza a vibrar (ruido de cuantización amplificado por las ganancias derivativas) o se vuelve inestable (demasiado retraso). Una alternativa más elaborada es un observador de Luenberger o un filtro de Kalman, que es el tema de la práctica~10.
\end{nota}

\chapter{Diseño de controladores}
\label{cap:control}

\section{Estructura general del lazo}

Todos los modos comparten la misma estructura (figura~\ref{fig:lazo}). La referencia $r$ pasa por un limitador de pendiente (\SI{1}{m/s}) y se satura a $\pm\SI{0.6}{m}$. El controlador discreto recibe solo las posiciones medidas, estima las velocidades y calcula $F_{\text{cmd}}$. A esa fuerza se le suman la fuerza manual y una señal externa opcional $u_{\text{ext}}$, y el resultado pasa por el actuador saturado.

\begin{figure}[H]
\centering
\begin{tikzpicture}[>=Latex,node distance=0.7cm and 0.75cm,
  b/.style={draw,thick,rounded corners=2pt,minimum height=1cm,minimum width=1.5cm,align=center,font=\small,fill=#1},
  b/.default=c1!8,
  s/.style={draw,circle,thick,inner sep=1.5pt}]
\node (r) {$r$};
\node[b=gray!10,right=of r] (ref) {Limitador\\pendiente\\+ sat.};
\node[b=c4!10,right=1.0cm of ref,minimum height=1.8cm] (ctrl) {Estimador\\$\hat{\vect s}$\\+\\Ley de\\control};
\node[s,right=of ctrl] (sum) {$+$};
\node[b=c2!10,right=of sum] (act) {$\dfrac{1}{\tau_a s+1}$\\ sat. $\pm F_{\max}$};
\node[b=c5!10,right=of act,minimum height=1.8cm] (pl) {Carro-\\péndulo\\no lineal};
\node[right=0.8cm of pl] (out) {$\vect s$};
\node[b=yellow!15,below=1.0cm of act] (sen) {Encoders\\$\mathcal Q_q$, ZOH};
\node[above=0.6cm of sum,font=\small] (uext) {$u_{\text{man}}+u_{\text{ext}}$};
\node[above=0.6cm of pl,font=\small] (fd) {$F_d$};
\draw[->] (r) -- (ref); \draw[->] (ref) -- node[above,font=\small]{$r'$} (ctrl);
\draw[->] (ctrl) -- node[above,font=\small]{$F_c$} (sum);
\draw[->] (uext) -- (sum);
\draw[->] (sum) -- node[above,font=\small]{$F_{\text{cmd}}$} (act); \draw[->] (act) -- node[above,font=\small]{$F$} (pl);
\draw[->] (fd) -- (pl);
\draw[->] (pl) -- (out);
\draw[->] ($(pl.east)+(0.4,0)$) |- (sen);
\draw[->] (sen) -| node[pos=0.25,below,font=\small]{$\vect y_k=[x_m,\ \theta_m]$} (ctrl);
\begin{scope}[on background layer]
\node[draw,dashed,c4,fit=(ctrl),inner sep=6pt,label={[c4,font=\footnotesize]above:discreto, $T_s=\SI{2}{ms}$}] {};
\end{scope}
\end{tikzpicture}
\caption{Lazo de control común a todos los modos.}
\label{fig:lazo}
\end{figure}

\section{Regulador lineal cuadrático (LQR)}
\label{sec:lqr}

\subsection{Planteamiento}

Dado el modelo lineal $\dot{\vect s}=A\vect s+Bu$, el LQR busca la realimentación $u=-K\vect s$ que minimiza
\begin{equation}
J = \int_0^\infty \left(\vect s^\top Q\,\vect s + u^\top R\,u\right)\dd t,\qquad Q=Q^\top\succeq0,\ R=R^\top\succ0 .
\end{equation}
Si $(A,B)$ es estabilizable y $(A,Q^{1/2})$ detectable, la solución es
\begin{equation}
K = R^{-1}B^\top X,
\end{equation}
donde $X=X^\top\succ0$ es la solución estabilizante de la \textbf{ecuación algebraica de Riccati}
\begin{equation}
A^\top X + XA - XBR^{-1}B^\top X + Q = 0 .
\label{eq:are}
\end{equation}
El lazo cerrado $A-BK$ es siempre estable y, para una sola entrada, garantiza un margen de ganancia de $[\tfrac12,\infty)$ y un margen de fase de al menos $60^\circ$, \emph{siempre que se realimente el estado real}. En el banco, la estimación de velocidades y el retraso del actuador reducen estos márgenes.

\subsection{Resolución por el método del Hamiltoniano}
\label{sec:hamiltoniano}

Cuando no se dispone de \code{lqr()}, \code{cp\_lqr.m} construye la matriz Hamiltoniana
\begin{equation}
\mathcal H=\begin{bmatrix} A & -BR^{-1}B^\top\\ -Q & -A^\top\end{bmatrix}\in\mathbb R^{2n\times 2n}.
\end{equation}
Sus autovalores son simétricos respecto al eje imaginario. Si las columnas de $\left[\begin{smallmatrix}V_1\\V_2\end{smallmatrix}\right]$ generan el subespacio invariante asociado a los $n$ autovalores estables, entonces $X=V_2V_1^{-1}$ y los autovalores estables de $\mathcal H$ son precisamente los polos del lazo cerrado.

\begin{lstlisting}[caption={Resolución de Riccati sin toolboxes (extracto de \code{cp\_lqr.m}).}]
H  = [A, -B/R*B'; -Q, -A'];
[V, Dg] = eig(H);
Vs = V(:, real(diag(Dg)) < 0);          % n vectores estables
X  = real(Vs(n+1:end,:) / Vs(1:n,:));
K  = R \ (B'*X);
\end{lstlisting}

\subsection{Elección de los pesos}

Una regla de partida útil es la \emph{regla de Bryson}: $Q_{ii}=1/s_{i,\max}^2$ y $R=1/u_{\max}^2$, donde $s_{i,\max}$ y $u_{\max}$ son los valores máximos aceptables de cada variable. Después se ajusta a la vista de los resultados:
\begin{itemize}
  \item aumentar $Q_{\theta}$ hace el lazo del ángulo más rápido y rígido, a costa de más fuerza;
  \item aumentar $Q_x$ hace que el carro vuelva antes a la referencia, pero con más inclinación transitoria;
  \item aumentar $R$ suaviza la fuerza, aleja el riesgo de saturación y vuelve el sistema más lento.
\end{itemize}

\begin{ejemplo}{Diseño por defecto}
Con $Q=\operatorname{diag}(50,1,200,1)$ y $R=0{,}1$:
\[
K_\uparrow = \begin{bmatrix}-22{,}36 & -19{,}55 & -84{,}23 & -14{,}27\end{bmatrix},\qquad
\lambda(A-BK)=\{-13{,}46\pm j5{,}55,\ -1{,}70\pm j1{,}35\}.
\]
Se observa que $K_x=-\sqrt{Q_x/R}=-\sqrt{500}$. Es una propiedad general del LQR cuando el estado tiene un modo integrador puro sin coste cruzado (aquí, la posición del carro). El par de polos rápidos corresponde al ángulo y el lento, al carro. Los signos negativos de $K$ se explican porque, con $u=-K\vect s$, una inclinación $\theta>0$ produce $F>0$: el carro se mueve para ponerse debajo de la varilla.
\end{ejemplo}

\begin{figure}[H]
\centering
\begin{tikzpicture}
\begin{groupplot}[group style={group size=1 by 3,vertical sep=0.35cm,xticklabels at=edge bottom},
  manual,height=3.6cm,xmin=0,xmax=5]
\nextgroupplot[ylabel={$\theta$ [$^\circ$]},legend pos=north east,legend columns=3]
\addplot[c1] table[x=t,y=thdeg]{data/qsweep_20.dat};
\addplot[c2] table[x=t,y=thdeg]{data/qsweep_200.dat};
\addplot[c4] table[x=t,y=thdeg]{data/qsweep_2000.dat};
\legend{$Q_\theta=20$,$Q_\theta=200$,$Q_\theta=2000$}
\nextgroupplot[ylabel={$x$ [m]}]
\addplot[c1] table[x=t,y=x]{data/qsweep_20.dat};
\addplot[c2] table[x=t,y=x]{data/qsweep_200.dat};
\addplot[c4] table[x=t,y=x]{data/qsweep_2000.dat};
\nextgroupplot[ylabel={$F$ [N]},xlabel={$t$ [s]}]
\addplot[c1] table[x=t,y=F]{data/qsweep_20.dat};
\addplot[c2] table[x=t,y=F]{data/qsweep_200.dat};
\addplot[c4] table[x=t,y=F]{data/qsweep_2000.dat};
\end{groupplot}
\end{tikzpicture}
\caption{Efecto de $Q_\theta$ partiendo de $\theta_0=\SI{0.15}{rad}$ ($8{,}6^\circ$). Con $Q_\theta=2000$ la fuerza llega a saturar a \SI{15}{N}.}
\label{fig:qsweep}
\end{figure}

\begin{figure}[H]
\centering
\begin{tikzpicture}
\begin{axis}[manual,width=0.8\linewidth,height=6cm,xlabel={$\mathrm{Re}(\lambda)$},ylabel={$\mathrm{Im}(\lambda)$},
  colorbar,colormap/viridis,point meta={log10(\thisrow{q})},
  colorbar style={ylabel={$\log_{10}Q_\theta$},width=0.25cm},xmax=0]
\addplot[scatter,only marks,mark=*,mark size=1.3pt] table[x=re,y=im]{data/lqr_locus.dat};
\end{axis}
\end{tikzpicture}
\caption{Lugar de los polos del lazo cerrado LQR al variar $Q_\theta$ entre $1$ y $10^4$ ($Q_x=50$, $R=0{,}1$). El par rápido, el del ángulo, se aleja del origen; el par lento, el del carro, apenas se mueve.}
\label{fig:locus}
\end{figure}

\subsection{Seguimiento de una referencia de posición}
\label{sec:seguimiento}

Como $x$ es un modo integrador, para cualquier $r$ constante el estado $\vect s_r=[r,0,0,0]^\top$ es un equilibrio. Basta, por tanto, con realimentar el error:
\begin{equation}
F = -K(\vect s-\vect s_r) = -K_x(x-r) - K_{\xd}\xd - K_\theta\theta - K_{\thd}\thd .
\end{equation}
En régimen permanente el error es nulo frente a escalones, salvo por la fricción de Coulomb, que deja un pequeño error. Frente a rampas o senos el carro va con retraso. La respuesta a un escalón (figura~\ref{fig:lqrstep}) muestra la \textbf{respuesta inversa} anunciada en la sección~\ref{sec:fnm}: el carro retrocede unos \SI{3}{cm} antes de avanzar.

\begin{figure}[H]
\centering
\begin{tikzpicture}
\begin{groupplot}[group style={group size=1 by 3,vertical sep=0.35cm,xticklabels at=edge bottom},
  manual,height=3.5cm,xmin=0,xmax=8]
\nextgroupplot[ylabel={$x$ [m]},legend pos=south east]
\addplot[c1] table[x=t,y=x]{data/lqr_step.dat};
\addplot[black,dashed,const plot] table[x=t,y=r]{data/lqr_step.dat};
\legend{$x$,$r$}
\draw[c2,thick,->] (axis cs:2.2,-0.12) -- (axis cs:1.25,-0.03) node[pos=0,right,font=\footnotesize]{respuesta inversa};
\nextgroupplot[ylabel={$\theta$ [$^\circ$]}]
\addplot[c2] table[x=t,y=thdeg]{data/lqr_step.dat};
\nextgroupplot[ylabel={$F$ [N]},xlabel={$t$ [s]}]
\addplot[c4] table[x=t,y=F]{data/lqr_step.dat};
\end{groupplot}
\end{tikzpicture}
\caption{LQR: escalón de \SI{0.3}{m} en la referencia, sin limitador de pendiente. Para avanzar, el carro primero inclina la varilla hacia delante, retrocediendo $\approx\SI{3}{cm}$. Con el limitador por defecto (\code{P.ref\_rate} $=\SI{1}{m/s}$) el retroceso se reduce a $\approx\SI{1.9}{cm}$ (apéndice~\ref{cap:validacion}).}
\label{fig:lqrstep}
\end{figure}

\begin{figure}[H]
\centering
\begin{tikzpicture}
\begin{groupplot}[group style={group size=2 by 1,horizontal sep=1.6cm},
  manual,width=0.5\linewidth,height=4.5cm,xmin=0,xmax=6,xlabel={$t$ [s]}]
\nextgroupplot[ylabel={$\theta$ [$^\circ$] / $x$ [cm]},legend pos=north east,title={\small Desde $\theta_0=0{,}15$ rad}]
\addplot[c2] table[x=t,y=thdeg]{data/lqr_ic.dat};
\addplot[c1] table[x=t,y expr=100*\thisrow{x}]{data/lqr_ic.dat};
\legend{$\theta$,$x$}
\nextgroupplot[ylabel={$\theta$ [$^\circ$] / $x$ [cm]},title={\small Empujón de \SI{1}{N}$\times$\SI{0.1}{s} en $t=1$ s}]
\addplot[c2] table[x=t,y=thdeg]{data/lqr_push.dat};
\addplot[c1] table[x=t,y expr=100*\thisrow{x}]{data/lqr_push.dat};
\end{groupplot}
\end{tikzpicture}
\caption{LQR: recuperación desde una condición inicial (asentamiento a $0{,}5^\circ$ en \SI{1.7}{s}, $F_{\max}=\SI{12.7}{N}$) y rechazo de un empujón en la punta (el carro recorre hasta \SI{37}{cm}).}
\label{fig:lqric}
\end{figure}

\section{Control PID del ángulo con PD del carro}
\label{sec:pid}

Es la estructura más intuitiva: un PID que mantiene la varilla vertical, más un lazo PD que devuelve el carro a la referencia:
\begin{equation}
F = \underbrace{K_{p\theta}\,\theta + K_{i\theta}\!\int\!\theta\,\dd t + K_{d\theta}\,\thd}_{\text{equilibrio}}
  \;+\; \underbrace{K_{px}(x-r) + K_{dx}\,\xd}_{\text{posición del carro}} .
\label{eq:pid}
\end{equation}

\begin{aviso}[title={¿Por qué $K_{px}>0$?}]
Parece contraintuitivo: si el carro está a la derecha de la referencia ($x-r>0$), el término de posición lo empuja \emph{todavía más} a la derecha. Es consecuencia de la fase no mínima. Para volver a la izquierda, la varilla tiene que inclinarse primero hacia la izquierda, y eso se consigue moviendo el carro a la derecha. Si se cambia el signo de $K_{px}$, el lazo se vuelve inestable.
\end{aviso}

Sin la parte integral, \eqref{eq:pid} es una realimentación de estado con $K_{\text{PID}}=-[K_{px},\ K_{dx},\ K_{p\theta},\ K_{d\theta}]$. Esto permite comparar directamente con el LQR. Con los valores por defecto, $(K_{px},K_{dx},K_{p\theta},K_{d\theta})=(10,10,60,8)$, los polos son $-8{,}56\pm j10{,}52$ y $-1{,}03\pm j1{,}16$: el lazo del ángulo es más oscilante y el del carro, algo más lento que con el LQR.

\begin{figure}[H]
\centering
\begin{tikzpicture}
\begin{groupplot}[group style={group size=2 by 1,horizontal sep=1.6cm},
  manual,width=0.5\linewidth,height=4.5cm,xlabel={$t$ [s]}]
\nextgroupplot[xmin=0,xmax=4,ylabel={$\theta$ [$^\circ$]},legend pos=north east,title={\small Condición inicial $\theta_0=0{,}1$ rad}]
\addplot[c2] table[x=t,y=thdeg]{data/pid_ic.dat};
\addplot[c1] table[x=t,y=thdeg]{data/lqr_ic.dat};
\legend{PID+PD,LQR ($\theta_0=0{,}15$)}
\nextgroupplot[xmin=0,xmax=10,ylabel={$x$ [m]},legend pos=south east,title={\small Escalón de \SI{0.3}{m}}]
\addplot[c2] table[x=t,y=x]{data/pid_step.dat};
\addplot[black,dashed] table[x=t,y=r]{data/pid_step.dat};
\legend{$x$,$r$}
\end{groupplot}
\end{tikzpicture}
\caption{Modo PID+PD. El escalón en posición muestra un sobreimpulso y una cola lenta, debidos al par de polos lentos poco amortiguado.}
\label{fig:pid}
\end{figure}

\section{\emph{Swing-up} por control de energía}
\label{sec:swingup}

Partiendo de la varilla colgando, un controlador lineal no puede levantarla: el ángulo recorre $\pi$ radianes y la fuerza está limitada. La estrategia de Åström y Furuta~\cite{astrom2000} consiste en controlar la \textbf{energía} $E$ de la varilla, definida en~\eqref{eq:energia}, hasta llevarla al valor que corresponde a la varilla arriba en reposo ($E=0$). Cuando la varilla pasa cerca de la vertical, se conmuta al LQR.

\subsection{Ley de control}

Por~\eqref{eq:Edot}, despreciando $b_p$, se tiene $\dot E=-ml\cos\theta\,\thd\,a$, donde $a=\xdd$ es la aceleración del carro. Tomemos como función de Lyapunov $V=\tfrac12E^2$ y elijamos
\begin{equation}
a = k_E\,E\,\thd\cos\theta \quad\Longrightarrow\quad
\dot V = E\dot E = -k_E\,ml\,E^2(\thd\cos\theta)^2 \le 0 .
\end{equation}
La energía converge a $0$ salvo en los instantes en que $\thd\cos\theta=0$, que son aislados. En la práctica se satura la aceleración y se añade un término que mantiene el carro centrado:
\begin{equation}
a = \operatorname{sat}_{a_{\max}}\!\Big(k_E\,E\,\thd\cos\theta \;-\; k_{x}(x-r) \;-\; k_v\xd\Big).
\label{eq:swing}
\end{equation}
Si la varilla está exactamente en reposo abajo, $\thd=0$ y la ley no hace nada. En ese caso se añade un pequeño impulso inicial ($0{,}3\,a_{\max}$) mientras $|\thd|<\SI{0.05}{rad/s}$ y $\cos\theta<-0{,}99$.

\subsection{De la aceleración a la fuerza (dinámica inversa)}

La ley~\eqref{eq:swing} da una aceleración deseada. La fuerza necesaria se obtiene de las ecuaciones del modelo: primero se calcula la aceleración angular resultante de~\eqref{eq:el2} y después la fuerza de~\eqref{eq:el1}:
\begin{align}
\thdd^\ast &= \frac{mgl\sin\theta - ml\cos\theta\,a - b_p\thd}{\Jp},\\
F &= (M+m)\,a + ml\cos\theta\,\thdd^\ast - ml\sin\theta\,\thd^2 + b_c\xd .
\end{align}

\subsection{Conmutación}

La figura~\ref{fig:fsm} muestra la máquina de estados. Se usa histéresis: se captura al bajar de $|\theta|<\SI{0.35}{rad}$ ($20^\circ$) y se libera al superar $|\theta|>\SI{0.8}{rad}$. Así se evitan conmutaciones repetidas en el borde de la región.

\begin{figure}[H]
\centering
\begin{tikzpicture}[>=Latex,node distance=5cm,
  st/.style={draw,thick,rounded corners=8pt,minimum width=3.2cm,minimum height=1.3cm,align=center,font=\small}]
\node[st,fill=c3!20] (sw) {\textbf{SWING-UP}\\$a$ según \eqref{eq:swing}};
\node[st,fill=c5!20,right=of sw] (lq) {\textbf{LQR}\\$F=-K_\uparrow(\vect s-\vect s_r)$};
\draw[->,thick] (sw.north east) to[bend left=25] node[above,font=\small]{$|\theta_w|<\theta_{\text{catch}}=0{,}35$} (lq.north west);
\draw[->,thick] (lq.south west) to[bend left=25] node[below,font=\small]{$|\theta_w|>\theta_{\text{release}}=0{,}8$} (sw.south east);
\draw[<-,thick] (sw.west) -- ++(-1,0) node[left,font=\small]{inicio};
\end{tikzpicture}
\caption{Máquina de estados del modo 2. $\theta_w$ es el ángulo envuelto a $[-\pi,\pi)$.}
\label{fig:fsm}
\end{figure}

\begin{figure}[H]
\centering
\begin{tikzpicture}
\begin{groupplot}[group style={group size=1 by 4,vertical sep=0.3cm,xticklabels at=edge bottom},
  manual,height=3.3cm,xmin=0,xmax=10]
\nextgroupplot[ylabel={$\theta$ [$^\circ$]},ytick={-180,-90,0,90,180}]
\addplot[c2,only marks,mark size=0.4pt] table[x=t,y=thdeg]{data/swingup.dat};
\nextgroupplot[ylabel={$E$ [J]}]
\addplot[c4] table[x=t,y=E]{data/swingup_E.dat};
\addplot[c5!70!black,dashed,const plot] table[x=t,y expr=\thisrow{caught}*0.2-1.3]{data/swingup.dat};
\node[font=\footnotesize,c5!60!black,anchor=west] at (axis cs:6.0,-1.0) {capturado (LQR)};
\nextgroupplot[ylabel={$x$ [m]}]
\addplot[c1] table[x=t,y=x]{data/swingup.dat};
\nextgroupplot[ylabel={$F$ [N]},xlabel={$t$ [s]}]
\addplot[c4!60!black] table[x=t,y=F]{data/swingup.dat};
\end{groupplot}
\end{tikzpicture}
\caption{\emph{Swing-up} desde la varilla colgando en reposo. La energía sube por bombeo desde $-2mgl=\SI{-1.18}{J}$; la captura se produce en $t\approx\SI{5.5}{s}$ y el carro no supera $|x|=\SI{0.47}{m}$.}
\label{fig:swingup}
\end{figure}

\section{Modo grúa: posicionamiento con anti-balanceo}
\label{sec:grua}

Con la varilla colgando, el problema es el de una grúa puente: llevar la carga a una posición sin que quede oscilando. Se diseña un LQR sobre el modelo linealizado en $\theta=\pi$, usando la desviación $\varphi=\operatorname{wrap}(\theta-\pi)$. Con $Q=\operatorname{diag}(50,1,20,1)$ y $R=0{,}1$:
\[
K_\downarrow=\begin{bmatrix}22{,}36 & 10{,}87 & 19{,}24 & 0{,}35\end{bmatrix},\qquad
\lambda=\{-11{,}72,\ -7{,}26,\ -1{,}69\pm j2{,}98\}.
\]
Si se anulan las ganancias del ángulo, el controlador se reduce a un PD sobre el carro, que lo posiciona igual de bien pero deja la carga oscilando. La comparación (figura~\ref{fig:grua}) es elocuente: \SI{0.34}{\degree} de balanceo residual con anti-balanceo frente a \SI{16.5}{\degree} sin él.

\begin{figure}[H]
\centering
\begin{tikzpicture}
\begin{groupplot}[group style={group size=1 by 2,vertical sep=0.35cm,xticklabels at=edge bottom},
  manual,height=4cm,xmin=0,xmax=8]
\nextgroupplot[ylabel={$x$ [m]},legend pos=south east]
\addplot[c1] table[x=t,y=x]{data/crane.dat};
\addplot[c2] table[x=t,y=x]{data/crane_nosway.dat};
\addplot[black,dashed] table[x=t,y=r]{data/crane.dat};
\legend{LQR anti-balanceo,solo PD carro,$r$}
\nextgroupplot[ylabel={$\theta-\pi$ [$^\circ$]},xlabel={$t$ [s]}]
\addplot[c1] table[x=t,y expr={mod(\thisrow{thdeg}+360,360)-180}]{data/crane.dat};
\addplot[c2] table[x=t,y expr={mod(\thisrow{thdeg}+360,360)-180}]{data/crane_nosway.dat};
\end{groupplot}
\end{tikzpicture}
\caption{Modo grúa: escalón de \SI{0.5}{m} con y sin realimentación del ángulo.}
\label{fig:grua}
\end{figure}

\section{Robustez frente a errores de modelo}

Las ganancias se calculan con los parámetros nominales. Para ver qué ocurre si la planta real es distinta, la figura~\ref{fig:robust} mantiene $K_\uparrow$ nominal (calculado con $M=\SI{0.5}{kg}$) y cambia la masa real del carro. El lazo sigue estable con masas reales de 0,3 a \SI{2}{kg} (de 0,6 a 4 veces la nominal); a medida que el carro se hace más pesado, el recorrido del carro aumenta.

\begin{figure}[H]
\centering
\begin{tikzpicture}
\begin{groupplot}[group style={group size=2 by 1,horizontal sep=1.6cm},
  manual,width=0.5\linewidth,height=4.5cm,xmin=0,xmax=6,xlabel={$t$ [s]}]
\nextgroupplot[ylabel={$\theta$ [$^\circ$]},legend pos=north east]
\addplot[c1] table[x=t,y=thdeg]{data/msweep_0.3.dat};
\addplot[c2] table[x=t,y=thdeg]{data/lqr_ic.dat};
\addplot[c4] table[x=t,y=thdeg]{data/msweep_1.0.dat};
\addplot[c5!70!black] table[x=t,y=thdeg]{data/msweep_2.0.dat};
\legend{$M=0{,}3$,$M=0{,}5$ (nom.),$M=1$,$M=2$}
\nextgroupplot[ylabel={$x$ [m]}]
\addplot[c1] table[x=t,y=x]{data/msweep_0.3.dat};
\addplot[c2] table[x=t,y=x]{data/lqr_ic.dat};
\addplot[c4] table[x=t,y=x]{data/msweep_1.0.dat};
\addplot[c5!70!black] table[x=t,y=x]{data/msweep_2.0.dat};
\end{groupplot}
\end{tikzpicture}
\caption{Robustez del LQR nominal frente a la masa real del carro.}
\label{fig:robust}
\end{figure}

\section{Supervisión y seguridad}

El controlador incluye un corte de seguridad con enclavamiento: si $|x_m|>x_{\text{safe}}=\SI{0.78}{m}$ en cualquier modo excepto el de lazo abierto, la fuerza de control se anula y se activa la señal \code{tripped}. Para rearmar el controlador hay que deshabilitarlo y volver a habilitarlo, o cambiar de modo. Este mecanismo protege frente a ganancias mal ajustadas y frente a perturbaciones grandes. Además, los topes elásticos del modelo de planta reproducen lo que ocurriría al llegar al final del raíl.

\section{Resumen comparativo}

\begin{table}[H]
\centering\small
\begin{tabularx}{\linewidth}{@{}lXXl@{}}
\toprule
Modo & Ventajas & Limitaciones & Parámetros\\
\midrule
1 LQR & óptimo, sistemático, buen compromiso entre fuerza y rapidez & necesita modelo y velocidades; válido solo cerca de $\theta=0$ & \code{Qup}, \code{Rup}\\
2 Swing-up & levanta la varilla desde cualquier estado & ajuste empírico; requiere recorrido de raíl & \code{k\_E}, \code{a\_max}, \code{kx\_sw}\\
3 PID+PD & intuitivo, sin modelo explícito & ajuste por prueba y error; signos poco intuitivos & \code{Kp\_th}…\code{Kd\_x}\\
4 Grúa & elimina el balanceo residual & solo con la varilla abajo & \code{Qdn}, \code{Rdn}\\
\bottomrule
\end{tabularx}
\caption{Comparación de los modos de control.}
\end{table}

\chapter{Implementación en Simulink}
\label{cap:simulink}

\section{Generación del modelo por script}

El modelo \code{cartpole\_sim.slx} no se dibuja a mano: lo genera \code{cp\_build\_model.m} mediante \code{add\_block} y \code{add\_line}. De esta forma el modelo es reproducible, se puede versionar con \code{git} como texto y se reconstruye en segundos tras cualquier cambio. El código de los bloques \emph{MATLAB Function} se lee de \code{cp\_plant\_block.m}, \code{cp\_controller\_block.m} y \code{cp\_ref\_block.m}, y se inserta mediante la API de Stateflow:

\begin{lstlisting}[caption={Inserción del código de un bloque \emph{MATLAB Function} y declaración de \code{P} como parámetro.}]
rt = sfroot;
ch = rt.find('-isa', 'Stateflow.EMChart', 'Path', blkPath);
ch.Script = fileread('cp_controller_block.m');
d = ch.find('-isa', 'Stateflow.Data', 'Name', 'P');
d.Scope = 'Parameter';         % P se lee del workspace (struct ajustable)
\end{lstlisting}

\begin{aviso}
Si se edita alguno de los archivos \code{*\_block.m}, hay que volver a ejecutar \code{cp\_build\_model}. Los cambios hechos directamente sobre el \code{.slx} se pierden al regenerarlo.
\end{aviso}

\section{Arquitectura}

\begin{figure}[H]
\centering
\resizebox{\linewidth}{!}{%
\begin{tikzpicture}[>=Latex,font=\small,
  ss/.style={draw,very thick,fill=#1,minimum width=2.2cm,minimum height=1.5cm,align=center,rounded corners=1pt},
  ss/.default=c1!10,
  cst/.style={draw,fill=yellow!25,minimum width=1.6cm,minimum height=0.6cm,font=\footnotesize\ttfamily},
  fw/.style={draw,fill=orange!20,minimum width=1.6cm,minimum height=0.6cm,font=\footnotesize\ttfamily},
  sm/.style={draw,circle,inner sep=1pt}]
\node[cst] (clk) at (0,2.2) {Clock};
\node[cst] (cref) at (0,1.4) {cmd\_ref};
\node[fw] (rext) at (0,0.6) {ref\_ext};
\node[ss=gray!12] (ref) at (2.8,1.4) {Reference};
\node[cst] (cctrl) at (2.8,-1.0) {cmd\_ctrl};
\node[ss=c4!12] (ctl) at (5.6,0.4) {Controller};
\node[sm] (sf) at (7.6,0.6) {$+$};
\node[fw] (uext) at (7.6,1.7) {u\_ext};
\node[ss=c2!12] (act) at (9.6,0.6) {Actuator};
\node[cst] (cd) at (11.4,2.3) {cmd\_dist};
\node[fw] (dext) at (12.9,2.3) {d\_ext};
\node[sm] (sd) at (12.2,1.4) {$+$};
\node[ss=c5!14] (pl) at (12.2,0.2) {Plant};
\node[ss=yellow!18] (sen) at (12.2,-2.0) {Sensors};
\node[draw,fill=gray!20,minimum height=4cm,minimum width=0.15cm] (mux) at (14.8,-0.4) {};
\node[draw,fill=c3!25,minimum width=1.5cm,minimum height=0.7cm,font=\footnotesize\ttfamily] (tel) at (16.4,-0.4) {telemetry};
\node[draw,fill=gray!10,minimum width=1.3cm,font=\footnotesize] (tw) at (18.4,0.3) {out.tel};
\node[draw,fill=gray!10,minimum width=1.3cm,font=\footnotesize] (sc) at (18.4,-1.1) {Scope};
\draw[->] (clk) -| ($(ref.west)+(-0.5,0.5)$) -- ($(ref.west)+(0,0.5)$);
\draw[->] (cref) -- (ref);
\draw[->] (rext) -| ($(ref.west)+(-0.5,-0.5)$) -- ($(ref.west)+(0,-0.5)$);
\draw[->] (ref.east) -- node[above]{$r$} ++(0.5,0) |- ($(ctl.west)+(0,0.3)$);
\draw[->] (cctrl) -| ($(ctl.west)+(-0.4,-0.45)$) -- ($(ctl.west)+(0,-0.45)$);
\draw[->] ($(ctl.east)+(0,0.2)$) -- node[above]{$F_c$} (sf);
\draw[->] (uext) -- (sf);
\draw[->] (sf) -- node[above]{$F_{\text{cmd}}$} (act);
\draw[->] (act) -- node[above]{$F$} ($(pl.west)+(0,0.4)$);
\draw[->] (cd) -| (sd); \draw[->] (dext) -| (sd);
\draw[->] (sd) -- node[right]{$F_d$} (pl);
\draw[->] (pl.east) -- node[above]{$\vect s$} ++(0.6,0) coordinate (b1) |- (sen.east);
\draw[->] (sen.west) -- ++(-1.5,0) -| node[pos=0.3,below]{$\vect y_m$} ($(ctl.south)+(-0.4,0)$);
\draw[->] (b1) -- ($(mux.west)+(0,0.6)$);
\draw[->,densely dotted] ($(ctl.east)+(0,-0.4)$) -- node[below,font=\footnotesize]{info} ++(0.7,0) |- (14.3,-3.2) |- ($(mux.west)+(0,-1.9)$);
\draw[->,densely dotted] (sen.east) ++(0,-0.3) -- ($(mux.west)+(0,-1.2)$);
\draw[->] (mux) -- (tel);
\draw[->] (tel) -- ++(0.9,0) |- (tw); \draw[->] (tel) -- ++(0.9,0) |- (sc);
\node[font=\scriptsize,align=left,anchor=north west] at (0,-1.7) {\colorbox{yellow!25}{\strut} Constant: modificable en vivo\\ \colorbox{orange!20}{\strut} From Workspace: se lee al arrancar};
\end{tikzpicture}}
\caption{Nivel superior del modelo \code{cartpole\_sim}. Los bloques amarillos son el «panel de mando» que modifican la clase \code{CartPole} y la GUI.}
\label{fig:arquitectura}
\end{figure}

\section{Subsistemas}

\begin{table}[H]
\centering\small
\begin{tabularx}{\linewidth}{@{}lp{3.2cm}X@{}}
\toprule
Subsistema & Entradas $\to$ salidas & Contenido\\
\midrule
\code{Reference} & $t$, \code{cmd\_ref}, \code{ref\_ext} $\to r$ & \emph{MATLAB Function} \code{cp\_ref\_block}: constante, seno, cuadrada, triangular o externa. Le siguen un \emph{Rate Limiter} ($\pm$\code{P.ref\_rate}) y una \emph{Saturation} ($\pm$\code{P.x\_ref\_max}).\\
\code{Controller} & $\vect y_m$, $r$, \code{cmd\_ctrl} $\to F_c$, info & ZOH sobre $r$ y \emph{MATLAB Function} \code{cp\_controller\_block}, con estado persistente (velocidades estimadas, integrador, captura, disparo). Se ejecuta a $T_s$.\\
\code{Actuator} & $F_{\text{cmd}}\to F$ & \emph{Transfer Fcn} $1/(\tau_as+1)$ más \emph{Saturation} $\pm F_{\max}$.\\
\code{Plant} & $F$, $F_d\to\vect s$ & \emph{MATLAB Function} \code{cp\_plant\_block} (modelo~\eqref{eq:nl}) más un \emph{Integrator} vectorial con condición inicial \code{P.x0}.\\
\code{Sensors} & $\vect s\to\vect y_m$ & Demux, ruido (\emph{Random Number}$\times\sigma\times$\code{noise\_on}), \emph{Quantizer} y ZOH a $T_s$.\\
\bottomrule
\end{tabularx}
\caption{Contenido de cada subsistema.}
\end{table}

\section{Configuración de la simulación}

\begin{itemize}
  \item \textbf{Solver}: ode4 (Runge--Kutta) de paso fijo, con $h=\SI{1}{ms}$. El paso fijo hace la simulación determinista y compatible con el \emph{pacing}. Los topes del raíl ($k_w=\SI{1e4}{N/m}$) tienen una frecuencia propia de unos \SI{140}{rad/s}, de modo que $\omega h\approx0{,}14$ y la integración sigue siendo estable.
  \item \textbf{Multifrecuencia}: la planta es continua y el controlador y los sensores son discretos a $T_s=2h$. Simulink inserta automáticamente las transiciones de frecuencia.
  \item \textbf{Parámetros}: todos los bloques leen sus parámetros del \emph{struct} \code{P} del workspace base. Con \code{DefaultParameterBehavior = Tunable}, basta con cambiar \code{P} y ejecutar\\ \code{set\_param(mdl, 'SimulationCommand', 'update')}\\ para aplicar el cambio en caliente.
  \item \textbf{Tiempo real}: \code{EnablePacing = on} y \code{PacingRate = 1} hacen que la simulación avance al ritmo del reloj de pared. Lo activa \code{CartPole.start()} y lo desactiva \code{CartPole.run()}.
  \item \textbf{Retrollamadas del modelo}: \code{InitFcn} crea \code{P} y las señales externas si no existen, para poder pulsar \emph{Run} directamente en Simulink. \code{StartFcn} y \code{StopFcn} conectan y desconectan el \emph{listener} de telemetría.
\end{itemize}

\section{Telemetría}

El bloque \code{telemetry} (una ganancia unidad, no virtual) concentra 13 señales. Esas señales se guardan en \code{out.tel} y además se leen en vivo con \code{add\_exec\_event\_listener}.

\begin{table}[H]
\centering\small
\begin{tabular}{@{}clll@{\hspace{1.2cm}}clll@{}}
\toprule
\# & Campo & Señal & Ud. & \# & Campo & Señal & Ud.\\
\midrule
1 & \code{x} & posición real & m & 8 & \code{Fact} & fuerza aplicada & N\\
2 & \code{xd} & velocidad real & m/s & 9 & \code{r} & referencia filtrada & m\\
3 & \code{th} & ángulo real & rad & 10 & \code{mode} & modo activo & --\\
4 & \code{thd} & vel. angular real & rad/s & 11 & \code{E} & energía (swing-up) & J\\
5 & \code{xm} & $x$ medida & m & 12 & \code{caught} & capturado por LQR & 0/1\\
6 & \code{thm} & $\theta$ medida & rad & 13 & \code{tripped} & corte de seguridad & 0/1\\
7 & \code{Fcmd} & fuerza pedida & N & & & & \\
\bottomrule
\end{tabular}
\caption{Vector de telemetría (\code{CartPole.TEL}).}
\label{tab:tel}
\end{table}

\chapter{Interfaces de usuario: código y GUI}
\label{cap:codigo}

\section{La clase \code{CartPole}}

\code{CartPole} encapsula todo el acceso al modelo. Permite trabajar de dos maneras (figura~\ref{fig:modos}):
\begin{itemize}
  \item \textbf{Por lotes} (\code{run(T)}): simula $T$ segundos tan rápido como puede y devuelve un \emph{struct} con todas las señales. Es la forma adecuada para diseñar, hacer barridos y generar figuras.
  \item \textbf{En vivo} (\code{start}/\code{stop}): simula indefinidamente en tiempo real. Los comandos tienen efecto inmediato y la telemetría llega a las funciones registradas como \emph{callback}.
\end{itemize}

\begin{figure}[H]
\centering
\resizebox{\linewidth}{!}{%
\begin{tikzpicture}[>=Latex,font=\small,
  ll/.style={draw,thick,minimum width=2.6cm,minimum height=0.8cm,fill=#1,align=center},ll/.default=gray!10,
  lab/.style={above,font=\scriptsize,fill=white,inner sep=1pt}]
\def\xa{0}\def\xb{6.5}\def\xc{13}
\node[ll=c1!12] (u) at (\xa,0) {Usuario /\\script / GUI};
\node[ll=c4!12] (c) at (\xb,0) {\code{CartPole}};
\node[ll=c5!12] (s) at (\xc,0) {Simulink\\\code{cartpole\_sim}};
\foreach \n in {u,c,s} \draw[dashed,gray] (\n.south) -- ++(0,-8.4);
\draw[->] (\xa,-1.1) -- node[lab]{\code{setMode}, \code{setRef}, \code{setInitialState}} (\xb,-1.1);
\draw[->] (\xb,-1.8) -- node[lab]{\code{set\_param(cmd\_*)}, \code{assignin('base','P')}} (\xc,-1.8);
\draw[->] (\xa,-2.6) -- node[lab]{\code{start()}} (\xb,-2.6);
\draw[->] (\xb,-3.3) -- node[lab]{\code{SimulationCommand=start} (pacing 1x)} (\xc,-3.3);
\draw[->,c2] (\xc,-4.0) -- node[lab]{\code{StartFcn}: \emph{listener} sobre \code{telemetry}} (\xb,-4.0);
\draw[->,c2] (\xc,-4.7) -- node[lab]{\code{PostOutputs} (cada paso)} (\xb,-4.7);
\draw[->,c2] (\xb,-5.4) -- node[lab]{\code{callback(t,d)} a 30 Hz} (\xa,-5.4);
\draw[->] (\xa,-6.2) -- node[lab]{\code{push(F,dur)}, \code{setParam(...)}} (\xb,-6.2);
\draw[->] (\xb,-6.9) -- node[lab]{\code{cmd\_dist} / \code{SimulationCommand=update}} (\xc,-6.9);
\draw[->] (\xa,-7.7) -- node[lab]{\code{stop(); res = lastResult()}} (\xb,-7.7);
\end{tikzpicture}}
\caption{Secuencia de interacción en el modo en vivo.}
\label{fig:modos}
\end{figure}

\subsection{Referencia de métodos}

\begin{table}[H]
\centering\small
\begin{tabularx}{\linewidth}{@{}p{5.6cm}Xc@{}}
\toprule
Método & Descripción & En vivo\\
\midrule
\code{cp = CartPole()} & Carga o genera el modelo, calcula $K$ y publica \code{P}. & --\\
\code{res = cp.run(T)} & Simulación por lotes de $T$ segundos. & no\\
\code{cp.start()}, \code{stop()}, \code{pause()}, \code{resume()} & Control de la simulación en vivo. & --\\
\code{cp.setMode(m)} & \code{'open'}, \code{'lqr'}, \code{'swingup'}, \code{'pid'}, \code{'crane'} (o 0..4). & sí\\
\code{cp.enable(tf)} & Habilita o inhibe el controlador (también rearma la seguridad). & sí\\
\code{cp.setForce(u)} & Fuerza manual en N, sumada a la del control. & sí\\
\code{cp.setRef(tipo,off,amp,f)} & Referencia: \code{'const'}, \code{'sine'}, \code{'square'}, \code{'triangle'}, \code{'external'}. & sí\\
\code{cp.push(F,dur)} & Empujón en la punta (temporizador). & sí\\
\code{cp.setParam(n,v)}, \code{setParams(S)} & Cambia parámetros, recalcula $K$ y lo aplica. & sí\\
\code{cp.setParamRaw(n,v)} & Cambia un parámetro sin recalcular $K$ (robustez, ganancias manuales). & sí\\
\code{cp.setInitialState(s0)} & Estado inicial para el próximo arranque. & no\\
\code{cp.setRefSignal(t,r)} & Referencia arbitraria (\code{ref\_ext}). & no\\
\code{cp.setForceSignal(t,u)} & Excitación de fuerza (\code{u\_ext}). & no\\
\code{cp.setDisturbanceSignal(t,Fd)} & Perturbación programada (\code{d\_ext}). & no\\
\code{cp.addTelemetryCallback(f)} & Registra \code{f(t,d)}, con \code{d} el vector de la tabla~\ref{tab:tel}. & --\\
\code{res = cp.lastResult()} & Datos de la última sesión en vivo. & --\\
\code{cp.animate(res, vel)} & Reproduce un ensayo. & --\\
\bottomrule
\end{tabularx}
\caption{API de \code{CartPole}.}
\end{table}

\subsection{Ejemplos}

\begin{ejemplo}{Ensayo por lotes y métricas}
\begin{lstlisting}
cp = CartPole();
cp.setInitialState([0 0 0.15 0]);
cp.setMode('lqr'); cp.setRef('const', 0);
res = cp.run(6);
ts  = res.t(find(abs(res.th_deg) > 0.5, 1, 'last'));   % tiempo de asentamiento
fprintf('ts = %.2f s,  Fmax = %.1f N\n', ts, max(abs(res.Fact)));
cp_plot_results(res, 'LQR');
\end{lstlisting}
\end{ejemplo}

\begin{ejemplo}{Sesión en vivo controlada desde un script}
\begin{lstlisting}
cp.setInitialState([0 0 0.05 0]); cp.setMode('lqr');
cp.TelemetryRate = 5;
cp.addTelemetryCallback(@(t,d) fprintf('%6.2f  x=%+.3f\n', t, d(CartPole.TEL.x)));
cp.start();               pause(3)       % la simulación corre en segundo plano
cp.setRef('sine', 0, 0.2, 0.2);  pause(5)
cp.push(1, 0.1);                  pause(3)
cp.stop();  res = cp.lastResult();  cp.clearCallbacks();
\end{lstlisting}
\end{ejemplo}

\begin{ejemplo}{Barrido de diseño: elegir $R$ para no saturar}
\begin{lstlisting}
Rs = logspace(-2, 1, 8);  Fmax = zeros(size(Rs));  ts = Fmax;
cp.setMode('lqr'); cp.setInitialState([0 0 0.15 0]);
for k = 1:numel(Rs)
    cp.setParam('Rup', Rs(k));            % recalcula K
    res = cp.run(6);
    Fmax(k) = max(abs(res.Fcmd));
    ts(k)   = res.t(find(abs(res.th_deg) > 0.5, 1, 'last'));
end
yyaxis left;  semilogx(Rs, Fmax, 'o-'); ylabel('F_{max} [N]')
yyaxis right; semilogx(Rs, ts, 's-');   ylabel('t_s [s]'); xlabel('R')
\end{lstlisting}
\end{ejemplo}

\subsection{Añadir un controlador propio}
\label{sec:nuevo_modo}

Para añadir, por ejemplo, un modo 5 (LQR con acción integral sobre el error de posición), hay que hacer tres cambios:
\begin{enumerate}
  \item En \code{cp\_params.m}, añadir \code{P.K\_int = zeros(1,5); P.Qint = diag([50 1 200 1 20]);}.
  \item En \code{cp\_design.m}, ampliar el modelo con el estado $\xi=\int(x-r)\,\dd t$ y calcular la ganancia:
\begin{lstlisting}
Aa = [A zeros(4,1); 1 0 0 0 0];   Ba = [B; 0];
P.K_int = cp_lqr(Aa, Ba, P.Qint, P.Rup);
\end{lstlisting}
  \item En \code{cp\_controller\_block.m}, añadir el caso:
\begin{lstlisting}
    case 5  % LQR con integral del error de posición
        I  = min(max(I + (x - r)*P.Ts, -P.I_max), P.I_max);
        Fc = -P.K_int*[x - r; xd; thw; thd; I];
\end{lstlisting}
\end{enumerate}
Por último, ejecutar \code{cp\_build\_model} y añadir \code{'lqri'} a \code{CartPole.MODES}. La práctica~4 propone evaluar este modo.

\section{Interfaz gráfica \code{cp\_gui}}

\code{cp\_gui} abre una ventana construida con \code{uifigure} (figura~\ref{fig:gui}). Usa internamente un objeto \code{CartPole}, de modo que todo lo que se hace en la GUI también puede hacerse desde la línea de comandos al mismo tiempo: \code{cp = cp\_gui;} devuelve el objeto.

\begin{figure}[H]
\centering
\begin{tikzpicture}[font=\footnotesize,x=1cm,y=1cm]
\draw[thick,fill=gray!5] (0,0) rectangle (15,9.4);
\fill[c1!70] (0,9.4) rectangle (15,9.9); \node[white,anchor=west] at (0.2,9.65) {Linear pendulum — live Simulink};
% animación
\draw[fill=white] (0.3,5.3) rectangle (10.3,9.2);
\node[anchor=north west,gray] at (0.35,9.15) {Animación};
\draw[thick,gray] (1.2,6.0) -- (9.4,6.0);
\draw[line width=2pt] (1.2,5.8) -- (1.2,6.3); \draw[line width=2pt] (9.4,5.8) -- (9.4,6.3);
\draw[fill=c1!60] (5.0,5.85) rectangle (5.9,6.3);
\draw[line width=3pt,c2] (5.45,6.2) -- (6.1,8.6); \fill[c2] (6.1,8.6) circle (0.1);
\fill[c5] (5.25,5.6) -- (5.45,5.85) -- (5.65,5.6) -- cycle;
% gráficas
\foreach \y/\l in {3.95/{$x$, $r$},2.6/{$\theta$},1.25/{$F$}} {
  \draw[fill=white] (0.3,\y-0.95) rectangle (10.3,\y+0.25);
  \node[anchor=north west,gray] at (0.35,\y+0.22) {\l\ (ventana de 10 s)};
}
\draw[c1,thick] plot[domain=0.6:10,samples=60] (\x,{3.55+0.25*sin(60*\x)});
\draw[c2,thick] plot[domain=0.6:10,samples=60] (\x,{2.2+0.2*exp(-0.3*\x)*sin(200*\x)});
\draw[c4,thick] plot[domain=0.6:10,samples=60] (\x,{0.85+0.25*exp(-0.3*\x)*cos(200*\x)});
% panel derecho
\begin{scope}[shift={(10.6,0)}]
\foreach \y/\t/\h in {7.6/{\textbf{Simulation}: \textcolor{c5!60!black}{$\blacktriangleright$Start} $\|$ \textcolor{c2}{$\blacksquare$Stop}\\ initial state, $\theta_0$, $x_0$, speed}/1.6,
                    5.5/{\textbf{Control}: mode $\blacktriangledown$, enabled $\square$\\ manual F (slider)\\ push}/1.8,
                    3.6/{\textbf{Cart reference}: type, offset,\\ amplitude, frequency}/1.6,
                    2.0/{\textbf{LQR tuning}: $Q_x, Q_{\dot x}, Q_\theta, R$\\ {[Recompute K]} \ $K=[\dots]$}/1.4,
                    0.2/{\textbf{Sensors and data}\\ \textbf{Status}: $t$, $x$, $\theta$, $F$, warnings}/1.6} {
  \draw[fill=white,rounded corners=2pt] (0,\y) rectangle (4.2,\y+\h);
  \node[anchor=north west,align=left,text width=4.05cm,font=\scriptsize] at (0.05,\y+\h-0.05) {\t};
}
\end{scope}
\end{tikzpicture}
\caption{Disposición de la interfaz gráfica.}
\label{fig:gui}
\end{figure}

\paragraph{Uso típico.}
\begin{enumerate}
  \item Elegir el estado inicial (arriba con $\theta_0$, o abajo) y la velocidad de simulación.
  \item Pulsar \textbf{Start}. Desde ese momento, cualquier cambio de modo, referencia, fuerza o pesos LQR se aplica en vivo.
  \item Con la ventana activa, las teclas $\leftarrow/\rightarrow$ aplican $\pm\SI{5}{N}$ mientras se mantienen pulsadas y la \textbf{barra espaciadora} da un empujón.
  \item Si aparece «SAFETY», hay que desmarcar y volver a marcar \emph{Control enabled}.
  \item \textbf{Stop} y después \textbf{Export last run}: los datos quedan en \code{cp\_res} y se representan con \code{cp\_plot\_results}.
\end{enumerate}

\begin{nota}
Prueba a levantar la varilla con el modo 0 y las flechas del teclado, sin controlador. Es muy difícil, y deja claro por qué el polo en $+5{,}45$ obliga a reaccionar en unas décimas de segundo. Después, con el modo 2, el controlador lo resuelve en unos 5 segundos.
\end{nota}

\chapter{Prácticas guiadas}
\label{cap:practicas}

Las prácticas están ordenadas de menor a mayor dificultad. Cada una indica sus objetivos, el procedimiento con el código necesario, las cuestiones a responder y lo que debe entregarse. Al empezar cada sesión conviene ejecutar:
\begin{lstlisting}
clear; close all; cp = CartPole(); cp.clearSignals();
\end{lstlisting}
Los valores numéricos que se citan como referencia corresponden a los parámetros por defecto.

% ----------------------------------------------------------------
\stepcounter{practica}
\begin{practica}{Exploración en lazo abierto y validación del modelo}
\objetivos
\begin{itemize}
  \item Familiarizarse con la interfaz y comprobar el modelo frente a la teoría.
  \item Medir la frecuencia y el amortiguamiento de la oscilación colgante.
\end{itemize}
\procedimiento
\begin{enumerate}
  \item Ensayo de caída libre desde arriba (sin control):
\begin{lstlisting}
cp.setMode('open'); cp.setInitialState([0 0 0.01 0]);
res = cp.run(3); cp_plot_results(res); cp.animate(res);
\end{lstlisting}
  \item Ajustar a los primeros \SI{0.5}{s} una exponencial $\theta(t)\approx\theta_0\cosh(pt)$ y estimar $p$.
  \item Oscilación libre colgando: \code{cp.setInitialState([0 0 pi-0.2 0]); res = cp.run(20);}. Medir el periodo y el decremento logarítmico $\delta=\ln(A_k/A_{k+1})$, y calcular $\zeta=\delta/\sqrt{4\pi^2+\delta^2}$.
  \item Repetir el punto anterior con el carro bloqueado (\code{cp.setParam('M',1e3)}) y comparar con $\omega_0=\sqrt{mgl/\Jp}$.
\end{enumerate}
\preguntas
\begin{enumerate}
  \item ¿Coincide $p$ con el polo inestable de la tabla~\ref{tab:polos} ($5{,}45$)? ¿Por qué no exactamente?
  \item Explicar por qué la frecuencia con el carro libre (\SI{0.89}{Hz}) es mayor que con el carro bloqueado (\SI{0.79}{Hz}).
  \item ¿Qué fricción domina el amortiguamiento de la oscilación, la del carro o la del pivote? Comprobarlo variando \code{bc} y \code{bp}.
\end{enumerate}
\entregables Gráficas de los tres ensayos y una tabla comparativa entre teoría y simulación.
\end{practica}

% ----------------------------------------------------------------
\stepcounter{practica}
\begin{practica}{Identificación experimental con una señal chirp}
\objetivos Obtener la respuesta en frecuencia del sistema colgante a partir de datos y compararla con la figura~\ref{fig:bode}.
\procedimiento
\begin{enumerate}
  \item Generar un chirp de fuerza de \SI{0.7}{N} entre \SI{0.5}{Hz} y \SI{3}{Hz} durante \SI{20}{s} y aplicarlo en lazo abierto:
\begin{lstlisting}
T = 20; tt = (0:cp.P.h:T)'; f0 = 0.5; f1 = 3;
u = 0.7*sin(2*pi*(f0*tt + (f1-f0)/(2*T)*tt.^2));
cp.setInitialState([0 0 pi 0]); cp.setMode('open');
cp.setForceSignal(tt, u); res = cp.run(T); cp.clearSignals();
\end{lstlisting}
  \item Estimar la respuesta en frecuencia con \code{tfestimate(res.Fact, mod(res.th,2*pi)-pi, [], [], [], 1/cp.P.h)}, de la Signal Processing Toolbox, o dividiendo directamente las FFT.
  \item Ajustar un modelo de segundo orden a $\Theta/F$ cerca de la resonancia.
\end{enumerate}
\preguntas
\begin{enumerate}
  \item ¿Dónde aparece la resonancia y dónde la antirresonancia de $X/F$?
  \item ¿Por qué no conviene excitar por debajo de \SI{0.3}{Hz} en lazo abierto? (Pista: polo en el origen y longitud del raíl.)
  \item ¿Cómo afecta la fricción de Coulomb a la identificación cuando la amplitud de la fuerza es pequeña?
\end{enumerate}
\entregables Bode experimental superpuesto al teórico y parámetros identificados.
\end{practica}

% ----------------------------------------------------------------
\stepcounter{practica}
\begin{practica}{Diseño de un LQR: influencia de $Q$ y $R$}
\objetivos Entender el compromiso entre rapidez, esfuerzo y recorrido del carro.
\procedimiento
\begin{enumerate}
  \item Desde $\theta_0=\SI{0.15}{rad}$, repetir el ensayo con $Q_\theta\in\{20,200,2000\}$ y luego con $R\in\{0{,}01,\ 0{,}1,\ 1\}$ (véase la figura~\ref{fig:qsweep}).
  \item Para cada diseño, anotar $K$, los polos del lazo cerrado (\code{cp.Lin.up.poles\_CL}), el tiempo de asentamiento a $0{,}5^\circ$, $F_{\max}$ y $|x|_{\max}$.
  \item Aplicar la regla de Bryson con $\theta_{\max}=5^\circ$, $x_{\max}=\SI{0.2}{m}$ y $F_{\max}=\SI{10}{N}$. Comparar el resultado con el diseño por defecto.
  \item En la GUI, cambiar $Q$ y $R$ en vivo mientras se aplican empujones.
\end{enumerate}
\preguntas
\begin{enumerate}
  \item Comprobar que $K_x=-\sqrt{Q_x/R}$ en todos los diseños y justificarlo.
  \item ¿Qué ocurre cuando el diseño satura el actuador? ¿Hay un ángulo inicial máximo recuperable? Encontrarlo por bisección.
\end{enumerate}
\entregables Tabla de diseños con sus métricas y una recomendación justificada.
\end{practica}

% ----------------------------------------------------------------
\stepcounter{practica}
\begin{practica}{Seguimiento de referencias y fase no mínima}
\objetivos Observar la respuesta inversa y cuantificar los límites de seguimiento.
\procedimiento
\begin{enumerate}
  \item Aplicar un escalón con \code{cp.setRefSignal([0 1 1.001 8], [0 0 0.3 0.3])} y medir el retroceso inicial (figura~\ref{fig:lqrstep}).
  \item Referencias senoidales de \SI{0.2}{m} a 0,05; 0,1; 0,2 y \SI{0.4}{Hz}: medir la ganancia y el desfase entre $x$ y $r$, y dibujar el Bode del lazo cerrado.
  \item Implementar el modo 5 con acción integral (sección~\ref{sec:nuevo_modo}) y repetir el ensayo de escalón con la fricción de Coulomb aumentada a \code{Fc = 1}.
\end{enumerate}
\preguntas
\begin{enumerate}
  \item Relacionar el ancho de banda medido con la posición del cero $z=4{,}91$.
  \item ¿Desaparece la respuesta inversa con el integrador? ¿Por qué no puede desaparecer con ningún controlador?
  \item Proponer un prefiltro o feedforward que reduzca el retraso con referencias senoidales.
\end{enumerate}
\entregables Gráficas, Bode del lazo cerrado y comparación con y sin integral.
\end{practica}

% ----------------------------------------------------------------
\stepcounter{practica}
\begin{practica}{Rechazo de perturbaciones y robustez}
\objetivos Evaluar el LQR frente a empujones y a errores en los parámetros.
\procedimiento
\begin{enumerate}
  \item Programar empujones de amplitud creciente (\SI{0.5}{N} a \SI{3}{N}, durante \SI{0.1}{s}) con \code{setDisturbanceSignal} y determinar la amplitud máxima admisible antes de que salte la seguridad.
  \item Sin recalcular $K$, variar la masa real del carro, la masa de la varilla y su longitud. Para ello, usar \code{cp.setParamRaw('M', 1)}, que cambia la planta sin recalcular $K$, en lugar de \code{setParam}, que sí lo recalcularía.
  \item Repetir con el ruido de los sensores activado (\code{cp.setParam('noise\_on',1)}).
\end{enumerate}
\preguntas
\begin{enumerate}
  \item ¿Cuál de los parámetros es más crítico? Relacionarlo con la expresión de $A$ en~\eqref{eq:AB}.
  \item ¿Qué margen de ganancia tiene el lazo real? Estimarlo multiplicando $K$ por un factor $\alpha$ hasta que el sistema se vuelva inestable, y comparar con el margen teórico del LQR.
\end{enumerate}
\entregables Mapa de estabilidad en función del parámetro y de $\alpha$.
\end{practica}

% ----------------------------------------------------------------
\stepcounter{practica}
\begin{practica}{PID+PD frente a LQR}
\objetivos Ajustar el controlador clásico y compararlo con el óptimo.
\procedimiento
\begin{enumerate}
  \item Con el modo \code{'pid'}, empezar con $K_{px}=K_{dx}=0$ y ajustar $K_{p\theta}$ y $K_{d\theta}$ hasta estabilizar el ángulo. Observar cómo deriva el carro.
  \item Añadir $K_{px}$ y $K_{dx}$, primero con signo negativo y luego con signo positivo. Explicar lo que ocurre.
  \item Calcular los polos de $A-BK_{\text{PID}}$ para cada ajuste y comparar con los del LQR.
  \item Activar $K_{i\theta}$ y aplicar una fuerza constante de \SI{1}{N} (\code{cp.setForce(1)}). ¿Qué hace el integrador?
\end{enumerate}
\preguntas
\begin{enumerate}
  \item ¿Puede un ajuste PID reproducir exactamente el LQR? Razonar sobre la estructura de~\eqref{eq:pid}.
  \item Ventajas prácticas del PID cuando el modelo es poco conocido.
\end{enumerate}
\entregables Ajuste final, polos y ensayo comparativo.
\end{practica}

% ----------------------------------------------------------------
\stepcounter{practica}
\begin{practica}{\emph{Swing-up} por energía}
\objetivos Entender el control no lineal basado en energía y la conmutación entre controladores.
\procedimiento
\begin{enumerate}
  \item Ejecutar el \emph{swing-up} por defecto (figura~\ref{fig:swingup}) y verificar que $E(t)$ crece de forma monótona hasta la captura.
  \item Barrer $k_E\in\{2,5,10,20\}$ y $a_{\max}\in\{2,3,4,6\}$. Registrar el tiempo de captura y $|x|_{\max}$, y marcar los casos que activan la seguridad.
  \item Eliminar el recentrado (\code{kx\_sw = 0}, \code{kv\_sw = 0}) y observar la deriva del carro.
  \item Reducir \code{th\_catch} a \SI{0.15}{rad} y aumentarlo a \SI{0.6}{rad}. ¿Qué limita la región de captura?
\end{enumerate}
\preguntas
\begin{enumerate}
  \item Demostrar que $\dot V\le0$ con la ley~\eqref{eq:swing} sin saturación ni recentrado.
  \item ¿Por qué no basta con un LQR con mucha ganancia para levantar la varilla?
  \item Proponer una forma de reducir el recorrido del carro, por ejemplo con una energía objetivo ligeramente mayor que cero o con bombeo en varias etapas.
\end{enumerate}
\entregables Mapa (tiempo de captura y recorrido) frente a $(k_E,a_{\max})$ y explicación de la conmutación.
\end{practica}

% ----------------------------------------------------------------
\stepcounter{practica}
\begin{practica}{Grúa con anti-balanceo}
\objetivos Diseñar un posicionamiento rápido sin oscilación residual de la carga.
\procedimiento
\begin{enumerate}
  \item Reproducir la figura~\ref{fig:grua}: el modo \code{'crane'} con referencia cuadrada (\code{cp.setRef('square',0,0.4,0.1)}), frente al mismo modo con las ganancias angulares a cero:
\begin{lstlisting}
Kn = cp.P.K_dn; Kn(3:4) = 0;
cp.setParamRaw('K_dn', Kn);        % sin recalcular el LQR
\end{lstlisting}
  \item Diseñar un \emph{input shaper} ZV (dos impulsos separados medio periodo de la oscilación) para la referencia del modo PD y compararlo con el LQR.
  \item Variar la longitud de la varilla $\pm30\,\%$ sin rediseñar. ¿Qué método es más robusto?
\end{enumerate}
\preguntas Comparar tiempo de posicionamiento, balanceo residual y sensibilidad a la longitud de la varilla.
\end{practica}

% ----------------------------------------------------------------
\stepcounter{practica}
\begin{practica}{Efecto de la instrumentación}
\objetivos Cuantificar cómo afectan el muestreo, la cuantización, el ruido y el filtro derivativo.
\procedimiento
\begin{enumerate}
  \item Barrer $T_s\in\{1,2,5,10,20\}$\,ms. Mantener $h\le T_s/2$ con \code{cp.setParam('h',min(1e-3,Ts/2))}; \code{run()} ajusta \code{FixedStep} automáticamente. ¿A partir de qué $T_s$ se degrada el lazo? ¿Y cuándo se vuelve inestable?
  \item Reducir la resolución del encoder angular a 1024, 512 y 256 cuentas por vuelta y observar el «ruido» en $F$.
  \item Barrer $\tau_v\in\{1,5,10,30,100\}$\,ms. Representar la varianza de $F$ en régimen permanente frente al tiempo de asentamiento.
\end{enumerate}
\preguntas
\begin{enumerate}
  \item Explicar el compromiso de $\tau_v$ (figura~\ref{fig:velest}).
  \item ¿Qué resolución mínima de encoder recomendaríais para el banco real?
\end{enumerate}
\end{practica}

% ----------------------------------------------------------------
\stepcounter{practica}
\begin{practica}{Proyecto abierto: observador de estado}
\objetivos Sustituir la derivada filtrada por un observador basado en el modelo.
\procedimiento
\begin{enumerate}
  \item Discretizar el modelo lineal: $A_d=e^{AT_s}$, $B_d=\int_0^{T_s}e^{A\tau}B\,\dd\tau$ (véase la sección~7 de \code{examples.m}).
  \item Diseñar un observador de Luenberger $\hat{\vect s}_{k+1}=A_d\hat{\vect s}_k+B_dF_k+L(\vect y_k-C\hat{\vect s}_k)$ con polos 3--5 veces más rápidos que el lazo, o un filtro de Kalman estacionario tomando como varianzas de medida $q^2/12+\sigma^2$.
  \item Implementarlo como modo 6 en \code{cp\_controller\_block.m}, con \code{persistent shat}, y comparar con el modo 1 con ruido activado.
  \item Probar a realimentar solo $x$ (observable) o solo $\theta$ (no observable).
\end{enumerate}
\entregables Código del modo nuevo, comparación de la varianza de $F$ y del tiempo de asentamiento, y discusión del principio de separación.
\end{practica}


\appendix
\chapter{Tabla de parámetros}
\label{ap:param}

\begin{table}[H]
\centering\footnotesize
\begin{tabularx}{\linewidth}{@{}p{4.1cm}p{3.4cm}p{2.1cm}X@{}}
\toprule
Campo de \code{P} & Valor & Unidad & Descripción\\
\midrule
\multicolumn{4}{@{}l}{\textit{Mecánica}}\\
\code{M}, \code{m} & 0,5 / 0,2 & kg & masa del carro / de la varilla\\
\code{Lp}, \code{l} & 0,6 / 0,3 & m & longitud de la varilla / distancia al CdM\\
\code{J} & 0,006 & kg\,m$^2$ & inercia de la varilla (CdM), $mL_p^2/12$\\
\code{bc}, \code{Fc}, \code{v\_eps} & 0,5 / 0,2 / 0,01 & N\,s/m, N, m/s & fricción viscosa y de Coulomb del carro\\
\code{bp} & 0,002 & N\,m\,s/rad & fricción del pivote\\
\code{x\_lim}, \code{k\_wall}, \code{b\_wall} & 0,8 / $10^4$ / 50 & m, N/m, N\,s/m & raíl y topes\\
\multicolumn{4}{@{}l}{\textit{Actuador y sensores}}\\
\code{Fmax}, \code{tau\_a} & 15 / 0,01 & N, s & saturación y dinámica del actuador\\
\code{res\_x}, \code{res\_th} & $10^{-4}$ / $2\pi/4096$ & m, rad & resolución de los encoders\\
\code{noise\_on}, \code{noise\_x}, \code{noise\_th} & 0 / $10^{-4}$ / $10^{-3}$ & --, m, rad & ruido de medida\\
\code{Ts}, \code{h}, \code{tau\_v} & 0,002 / 0,001 / 0,01 & s & muestreo, paso del solver, filtro derivativo\\
\multicolumn{4}{@{}l}{\textit{Referencia y seguridad}}\\
\code{x\_ref\_max}, \code{ref\_rate} & 0,6 / 1 & m, m/s & saturación y pendiente de la referencia\\
\code{x\_safe} & 0,78 & m & umbral del corte de seguridad\\
\code{x0} & $[0,0,0{,}1,0]^\top$ & -- & estado inicial\\
\multicolumn{4}{@{}l}{\textit{Controladores}}\\
\code{Qup}, \code{Rup} & diag(50,1,200,1) / 0,1 & -- & LQR arriba\\
\code{Qdn}, \code{Rdn} & diag(50,1,20,1) / 0,1 & -- & LQR grúa\\
\code{k\_E}, \code{a\_max} & 5 / 3 & --, m/s$^2$ & \emph{swing-up}\\
\code{kx\_sw}, \code{kv\_sw} & 8 / 1 & -- & recentrado durante el \emph{swing-up}\\
\code{th\_catch}, \code{th\_release} & 0,35 / 0,8 & rad & conmutación con histéresis\\
\code{Kp\_th}, \code{Ki\_th}, \code{Kd\_th} & 60 / 0 / 8 & -- & PID del ángulo\\
\code{Kp\_x}, \code{Kd\_x}, \code{I\_max} & 10 / 10 / 1 & -- & PD del carro y anti-windup\\
\bottomrule
\end{tabularx}
\caption{Parámetros por defecto (\code{cp\_params.m}).}
\end{table}

\chapter{Validación del modelo Simulink}
\label{cap:validacion}

Las gráficas de resultados de este manual se generaron con la réplica en Python del modelo (\code{python\_reference\_sim.py}). Para comprobar que representan fielmente al modelo Simulink, los mismos escenarios se ejecutaron en \code{cartpole\_sim.slx} (MATLAB R2025b, ode4, $h=\SI{1}{ms}$) y se compararon muestra a muestra. La réplica no modela el limitador de pendiente de la referencia, así que en los escenarios con escalón se desactivó (\code{cp.setParamRaw('ref\_rate', 1e6)}).

\begin{table}[H]
\centering\small
\begin{tabularx}{\linewidth}{@{}lrrX@{}}
\toprule
Escenario & RMS $\Delta x$ [mm] & RMS $\Delta\theta$ [$^\circ$] & Métrica (Python / Simulink)\\
\midrule
LQR desde $\theta_0=\SI{0.15}{rad}$ & 0,12 & 0,013 & asentamiento a $0{,}5^\circ$: 1,67 / 1,67 s; $F_{\max}$ 12,8 / 12,8 N\\
LQR, escalón \SI{0.3}{m} & 0,32 & 0,018 & retroceso inicial 29,4 / 29,4 mm\\
LQR, empujón \SI{1}{N}$\times$\SI{0.1}{s} & 1,34 & 0,044 & asentamiento 3,47 / 3,46 s\\
Swing-up + LQR & 0,41 & 0,19 & captura en 5,50 / 5,50 s; $|x|_{\max}$ 0,470 / 0,470 m\\
PID desde $\theta_0=\SI{0.1}{rad}$ & 0,11 & 0,010 & asentamiento 0,98 / 0,98 s\\
PID, escalón \SI{0.3}{m} & 0,12 & 0,006 & retroceso inicial 20,3 / 20,3 mm\\
Grúa con anti-balanceo & 0,17 & 0,016 & balanceo residual ($t>\SI{5}{s}$) 0,18 / 0,18$^\circ$\\
Grúa sin anti-balanceo & 0,17 & 0,057 & balanceo residual 15,25 / 15,24$^\circ$\\
\bottomrule
\end{tabularx}
\caption{Comparación entre la réplica en Python y el modelo Simulink.}
\label{tab:validacion}
\end{table}

Las discrepancias son del orden de la resolución de los encoders, por lo que las figuras del manual son representativas del modelo Simulink. Otros resultados obtenidos directamente en Simulink:
\begin{itemize}
  \item \textbf{Escalón con el limitador por defecto} (\SI{1}{m/s}): el retroceso inicial baja de 29 a 19~mm y la fuerza máxima de 5,5 a 2,1~N.
  \item \textbf{Ciclo límite}: tras el escalón, el carro oscila lentamente en torno a la referencia ($\approx\pm\SI{8}{mm}$ con LQR y $\approx\pm\SI{20}{mm}$ con PID) por la fricción de Coulomb, que ninguno de los dos controladores compensa al no tener acción integral sobre $x$.
  \item \textbf{Seguimiento de un seno de \SI{0.2}{m} con LQR}: retraso de fase de $28^\circ$ a \SI{0.1}{Hz}, $61^\circ$ a \SI{0.2}{Hz} y $140^\circ$ a \SI{0.5}{Hz} (ganancia 0,57).
  \item \textbf{Ruido en encoders activado}: desviación típica de $\theta$ en régimen de $0{,}08^\circ$ y de la fuerza de \SI{0.41}{N}.
  \item \textbf{Modo en vivo}: la simulación sigue el tiempo real tras la compilación inicial; los cambios de referencia, empujones y pesos LQR se aplican durante la ejecución.
\end{itemize}

Para repetir la validación:
\begin{lstlisting}
cd src/linear_pendulum/validation
run_validation                         % MATLAB: escribe results/*.csv
!python3 compare_with_python.py        % compara con la réplica en Python
\end{lstlisting}

\chapter{Resolución de problemas}

\begin{table}[H]
\centering\small
\begin{tabularx}{\linewidth}{@{}p{4.8cm}X@{}}
\toprule
Síntoma & Causa probable y solución\\
\midrule
\code{cp\_build\_model} falla en \code{setScript} & Versión de Stateflow sin la propiedad \code{Script}: actualizar a R2020b o posterior.\\
\emph{Undefined variable P} al pulsar \emph{Run} en Simulink & Ejecutar \code{CartPole()} o \code{P = cp\_design(cp\_params())}. La \code{InitFcn} lo hace sola si el modelo se ha generado con el script.\\
La simulación en vivo va más rápida que el tiempo real & La versión no soporta \emph{pacing} (anterior a R2018a). Reducir la carga o usar \code{run()}.\\
Error de tamaño al cambiar \code{P} en vivo & Los parámetros ajustables no pueden cambiar de dimensión durante la simulación. Hay que parar, cambiar y volver a arrancar.\\
El LQR no estabiliza y aparece «SEGURIDAD» & Ángulo inicial demasiado grande, empujón excesivo o $K$ mal calculado. Revisar \code{cp.P.K\_up} y deshabilitar y rehabilitar el control.\\
El \emph{swing-up} no captura y el carro golpea los topes & Bajar \code{a\_max} o \code{k\_E}, o subir \code{kx\_sw}.\\
\code{ref\_ext}, \code{u\_ext} o \code{d\_ext} no cambian en vivo & Los bloques \emph{From Workspace} leen la variable al arrancar. Para cambios en vivo, usar \code{setRef}, \code{setForce} o \code{push}.\\
La GUI se congela con \emph{PacingRate} alto & Bajar \code{TelemetryRate} o la velocidad de simulación.\\
\bottomrule
\end{tabularx}
\caption{Problemas habituales.}
\end{table}

\chapter{Resumen de fórmulas}

\begin{align*}
&\text{Modelo:} && \begin{bmatrix} M+m & ml\cos\theta\\ ml\cos\theta & \Jp\end{bmatrix}\begin{bmatrix}\xdd\\ \thdd\end{bmatrix}=\begin{bmatrix} F - F_f + ml\sin\theta\,\thd^2\\ mgl\sin\theta - b_p\thd\end{bmatrix}\\
&\text{Polo inestable (sin fricción):} && p=\sqrt{\frac{(M+m)mgl}{(M+m)\Jp-m^2l^2}}\\
&\text{Cero de fase no mínima de } X/F: && z=\sqrt{mgl/\Jp}\\
&\text{Energía:} && E=\tfrac12\Jp\thd^2+mgl(\cos\theta-1),\qquad \dot E=-ml\cos\theta\,\thd\,\xdd-b_p\thd^2\\
&\text{Swing-up:} && \xdd^\ast=\operatorname{sat}\big(k_EE\thd\cos\theta-k_x(x-r)-k_v\xd\big)\\
&\text{LQR:} && K=R^{-1}B^\top X,\quad A^\top X+XA-XBR^{-1}B^\top X+Q=0\\
&\text{Derivada filtrada:} && \hat v_k=a\hat v_{k-1}+(1-a)(y_k-y_{k-1})/T_s,\quad a=e^{-T_s/\tau_v}
\end{align*}

\begin{thebibliography}{9}
\bibitem{astrom2000} K.~J.~Åström y K.~Furuta, «Swinging up a pendulum by energy control», \emph{Automatica}, 36(2), pp.~287--295, 2000.
\bibitem{ogata} K.~Ogata, \emph{Ingeniería de control moderna}, 5.ª ed., Pearson, 2010.
\bibitem{franklin} G.~F.~Franklin, J.~D.~Powell y A.~Emami-Naeini, \emph{Feedback Control of Dynamic Systems}, 8.ª ed., Pearson, 2019.
\bibitem{anderson} B.~D.~O.~Anderson y J.~B.~Moore, \emph{Optimal Control: Linear Quadratic Methods}, Prentice Hall, 1990.
\bibitem{astromfb} K.~J.~Åström y R.~M.~Murray, \emph{Feedback Systems: An Introduction for Scientists and Engineers}, 2.ª ed., Princeton University Press, 2021.
\bibitem{skogestad} S.~Skogestad e I.~Postlethwaite, \emph{Multivariable Feedback Control: Analysis and Design}, 2.ª ed., Wiley, 2005 (limitaciones por ceros en el semiplano derecho).
\bibitem{singer} N.~C.~Singer y W.~P.~Seering, «Preshaping command inputs to reduce system vibration», \emph{ASME J. Dyn. Sys. Meas. Control}, 112(1), pp.~76--82, 1990.
\end{thebibliography}


\end{document}
